% BEGIN EMBEDDED BIBLIOGRAPHY
% Bibliography data are part of this source; the compiler writes its temporary input.
\begin{filecontents*}[overwrite]{\jobname.bib}
@unpublished{ChamberLattice,
 title={Reflection positivity for chamber {QED} on the infinite lattice},
 shorthand={CL}, date={2026}, note={Companion manuscript, version of 20 October 2026}}

@unpublished{ContinuumLimits,
 title={Gaussian domination and structural properties of continuum limits},
 shorthand={CE}, date={2026}, note={Companion manuscript, version of 20 October 2026}}

@unpublished{ErgodicComponents,
 title={Ergodic components and clustering for reflection-positive fields},
 shorthand={EC}, date={2026}, note={Companion manuscript, version of 20 October 2026}}

@unpublished{MaxwellBlock,
 title={Maxwell block minimization and Gaussian fluctuations},
 shorthand={MB}, date={2026}, note={Companion manuscript, version of 20 October 2026}}

@unpublished{FermionBlock,
 title={Chamber fermion locality and local integration},
 shorthand={FB}, date={2026}, note={Companion manuscript, version of 20 October 2026}}

@unpublished{NormalizedContraction,
 title={Relevant extraction and contraction of local interactions},
 shorthand={NC}, date={2026}, note={Companion manuscript, version of 20 October 2026}}

@unpublished{NonlinearStep,
 title={The exact nonlinear renormalization step for lattice {QED} in three dimensions},
 shorthand={NL}, date={2026}, note={Companion manuscript, version of 20 October 2026}}

@unpublished{CoverTuning,
 title={Agreement across covers and common tuning for lattice {QED} in three dimensions},
 shorthand={CT}, date={2026}, note={Companion manuscript, version of 20 October 2026}}

@unpublished{TerminalPressure,
 title={The terminal measure and pressure estimates for lattice {QED} in three dimensions},
 shorthand={TP}, date={2026}, note={Companion manuscript, version of 20 October 2026}}

@unpublished{PositiveProbe,
 title={Positive compact probes and free response coefficients for lattice {QED} in three dimensions},
 shorthand={PP}, date={2026}, note={Companion manuscript, version of 20 October 2026}}

@unpublished{ReferenceResponse,
 title={Exact source identities and normalized reference response for lattice {QED} in three dimensions},
 shorthand={RR}, date={2026}, note={Companion manuscript, version of 20 October 2026}}

@article{OS1973,
 author={Osterwalder, K. and Schrader, R.},
 title={Axioms for {Euclidean} {Green's} functions},
 journaltitle={Communications in Mathematical Physics},
 volume={31}, date={1973}, pages={83--112}, doi={10.1007/BF01645738}}
@article{OS1975,
 author={Osterwalder, K. and Schrader, R.},
 title={Axioms for {Euclidean} {Green's} functions {II}},
 journaltitle={Communications in Mathematical Physics},
 volume={42}, date={1975}, pages={281--305}, doi={10.1007/BF01608978}}
@article{OsterwalderSeiler1978,
 author={Osterwalder, K. and Seiler, E.},
 title={Gauge field theories on a lattice},
 journaltitle={Annals of Physics}, volume={110}, date={1978},
 pages={440--471}, doi={10.1016/0003-4916(78)90039-8}}

@article{FILS1978,
  author = {Fr\"ohlich, J. and Israel, R. and Lieb, E. H. and Simon, B.},
  title = {Phase transitions and reflection positivity.
    {I}. General theory and long range lattice models},
  journaltitle = {Communications in Mathematical Physics},
  volume = {62}, date = {1978}, pages = {1--34},
  doi = {10.1007/BF01940327}
}

@article{DimockUVS2022,
  author = {Dimock, J.},
  title = {Ultraviolet stability for {QED} in {$d=3$}},
  journal = {Annales Henri Poincar\'e},
  volume = {23},
  date = {2022},
  pages = {2113--2205},
  doi = {10.1007/s00023-021-01127-z},
  eprint = {2009.01156},
  eprinttype = {arxiv}
}

@unpublished{VolumeTransfer,
 title={Volume transfer and component cumulants for lattice {QED} in three dimensions},
 shorthand={VT}, date={2026}, note={Companion manuscript, version of 20 October 2026}}
\end{filecontents*}
% END EMBEDDED BIBLIOGRAPHY
% QEM1: modular revision of the existing field-strength construction.
\documentclass[11pt,letterpaper,reqno]{amsart}
% BEGIN INLINE manuscript-style
\makeatletter
\NeedsTeXFormat{LaTeX2e}[2020/10/01]
% Use with \documentclass[11pt,letterpaper,reqno]{amsart}.
% This file controls presentation only. It supplies no mathematical content.

\RequirePackage[T1]{fontenc}
\RequirePackage{lmodern}
\RequirePackage{microtype}
\RequirePackage[letterpaper,margin=1in,headheight=15pt]{geometry}
\RequirePackage{amsmath,amsthm,amssymb,mathtools}
\RequirePackage{enumitem}
\RequirePackage{perpage}
\RequirePackage{xcolor}
\definecolor{manuscriptred}{rgb}{0.6,0.1,0.1}
\definecolor{manuscriptblue}{rgb}{0.1,0.3,0.5}
\definecolor{manuscriptmagenta}{HTML}{60445C}
\RequirePackage[breaklinks=true,colorlinks=true,
  linkcolor=manuscriptred,citecolor=manuscriptblue,urlcolor=manuscriptmagenta]{hyperref}
\RequirePackage{cleveref}
\hypersetup{pdfauthor={}}

% Retain the AMS paragraph, title, heading and running-head conventions.
% Equations retain one global Arabic sequence, as in the reference.
\setcounter{tocdepth}{1}
\renewcommand{\thefigure}{\Roman{figure}}
\renewcommand{\thetable}{\Roman{table}}
\MakePerPage{footnote}
\renewcommand{\thefootnote}{%
  \ifnum\value{footnote}>9\arabic{footnote}\else\fnsymbol{footnote}\fi}

% Upright statements. Main results are A, B, C; supporting statements
% share a section-local sequence. These end marks indicate typography,
% never proof completeness, validation, or a deferred-proof status.
\theoremstyle{definition}
\newtheorem{ms@theorem}{Theorem}
\renewcommand{\thems@theorem}{\Alph{ms@theorem}}
\newtheorem*{ms@theoremstar}{Theorem}
\newtheorem{ms@proposition}{Proposition}[section]
\newtheorem{ms@lemma}[ms@proposition]{Lemma}
\newtheorem{ms@definition}[ms@proposition]{Definition}
\newtheorem{ms@corollary}[ms@proposition]{Corollary}
\newtheorem*{ms@corollarystar}{Corollary}

\newcommand{\ms@majorstart}{%
  \pushQED{\qed}\renewcommand{\qedsymbol}{$\blacksquare\square$}}
\newcommand{\ms@minorstart}{%
  \pushQED{\qed}\renewcommand{\qedsymbol}{$\blacksquare$}}
\newenvironment{theorem}[1][]{\ms@majorstart\ms@theorem[#1]}
  {\popQED\endms@theorem}
\newenvironment{theorem*}[1][]{\ms@majorstart\ms@theoremstar[#1]}
  {\popQED\endms@theoremstar}
\newenvironment{proposition}[1][]{\ms@minorstart\ms@proposition[#1]}
  {\popQED\endms@proposition}
\newenvironment{lemma}[1][]{\ms@minorstart\ms@lemma[#1]}
  {\popQED\endms@lemma}
\newenvironment{definition}[1][]{\ms@minorstart\ms@definition[#1]}
  {\popQED\endms@definition}
\newenvironment{corollary}[1][]{\ms@majorstart\ms@corollary[#1]}
  {\popQED\endms@corollary}
\newenvironment{corollary*}[1][]{\ms@majorstart\ms@corollarystar[#1]}
  {\popQED\endms@corollarystar}
% The source's lemma/proposition variants are deliberately not carried over:
% choose one marker for each statement type rather than encode proof status.

\theoremstyle{remark}
\newtheorem{remark}[ms@proposition]{Remark}
\newtheorem*{remark*}{Remark}
% The ordinary amsthm proof environment retains its open square.

\crefname{ms@theorem}{theorem}{theorems}
\Crefname{ms@theorem}{Theorem}{Theorems}
\crefname{ms@proposition}{proposition}{propositions}
\Crefname{ms@proposition}{Proposition}{Propositions}
\crefname{ms@lemma}{lemma}{lemmas}
\Crefname{ms@lemma}{Lemma}{Lemmas}
\crefname{ms@definition}{definition}{definitions}
\Crefname{ms@definition}{Definition}{Definitions}
\crefname{ms@corollary}{corollary}{corollaries}
\Crefname{ms@corollary}{Corollary}{Corollaries}
\crefname{remark}{remark}{remarks}
\Crefname{remark}{Remark}{Remarks}

% Bibliography configuration for this manuscript.
  \RequirePackage[backend=biber,style=alphabetic,backref=true,
    giveninits=true,isbn=false,date=year]{biblatex}
  \addbibresource{\jobname.bib}
  \setcounter{biburlnumpenalty}{100}
  \setcounter{biburllcpenalty}{100}
  \setcounter{biburlucpenalty}{100}
  \renewcommand*{\bibfont}{\footnotesize}
  \renewbibmacro{in:}{}
  \DeclareFieldFormat{pages}{#1}
  % Use the public alphabetic style. No private author-hash overrides;
  % repeated author names are printed in full.
\makeatother
% END INLINE manuscript-style
\usepackage{mathrsfs}
\allowdisplaybreaks
\newcommand{\R}{\mathbb R}
\newcommand{\C}{\mathbb C}
\newcommand{\Z}{\mathbb Z}
\newcommand{\N}{\mathbb N}
\newcommand{\E}{\mathbb E}
\newcommand{\Sspace}{\mathcal S}
\newcommand{\cS}{\mathcal S}
\newcommand{\dd}{\mathrm d}
\DeclareMathOperator{\supp}{supp}
\DeclareMathOperator{\tr}{tr}
\DeclareMathOperator{\diag}{diag}

\newcommand{\extref}[2]{#1 of~\cite{#2}}
\newcommand{\exteq}[2]{(#1) of~\cite{#2}}
\title[The continuum field-strength construction for QED3]{The continuum field-strength construction for QED in three dimensions}
\date{}
\hypersetup{pdftitle={The continuum field-strength construction for QED in three dimensions},pdfauthor={}}
\begin{document}
\begin{abstract}
At sufficiently small positive physical charge, almost every translation-ergodic component of every volume-first continuum limit of noncompact lattice QED with chamber fermions is a non-Gaussian law of closed two-form fields. Its Schwinger functions satisfy the Osterwalder--Schrader axioms, including clustering. The bare mass is chosen independently of the volume. For each compact probe with a positive free fourth coefficient, its fourth cumulant is positive at sufficiently small charge.
\end{abstract}
\maketitle
\tableofcontents
\section{Introduction}

The observable in this paper is the field strength. We regard it as a real
tempered two-form on $\R^3$, and write $F(f)$ for its pairing with a test
two-form $f$. The regularization uses a noncompact gauge field and
four-component chamber fermions, with all flux sectors included.
The order of limits is part of the construction: the volume tends to
infinity at a fixed lattice spacing, and the spacing then tends to zero.

For a lattice curvature $u$, the physical field $F^\sharp$ is the
centered plaquette step interpolation of $\eta^{-3/2}u$ at spacing
$\eta$; see \cref{sec:regularization,sec:continuum-limits} for the
finite law and interpolation. A volume-first limit is a weak limit
of fixed-spacing infinite-volume laws as $\eta\to0$. We take these
weak limits on the Hilbert space $H_{-4}$ of tempered two-forms
defined in \cref{sec:probability}. In the following statement
$\gamma_j=\operatorname{diag}(\sigma_j,-\sigma_j)$, where $\sigma_j$
are the Pauli matrices, and $d^*$ is the Euclidean adjoint of $d$.

\begin{theorem}[Field-strength OS theorem]\label{thm:main}
Fix $0<\rho<1$ and $0<\bar m\le1$, the chamber Wilson parameter and
reference mass parameter in physical units. Use the standard chamber weights
$(6,6\sqrt2,6\sqrt2,6)$. There are an odd scale factor $L\ge5$,
numbers $e_{\mathrm{phys}},\sigma>0$, and real volume-independent bare mass shifts
$m_{0,N}(e)$ for $N\in\Z_{\ge0}$ and $0<e\le e_{\mathrm{phys}}$, with
$|m_{0,N}(e)|<1/8$, with the following property.

Fix a real $\beta\in C_c^\infty(\R^3;\Lambda^2)$ such that
$f=dd^*\beta$ has
\begin{equation}\label{eq:positive-coefficient}
 b(f):=-6\tr\bigl[(S_{\bar m}\,i\gamma\cdot d^*\beta)^4\bigr]>0,
 \qquad S_{\bar m}=(\bar m+\gamma\cdot\partial)^{-1}.
\end{equation}
Such probes exist. There is $e_\beta\in(0,e_{\mathrm{phys}}]$ such that, for $0<e\le e_\beta$,
the all-flux chamber laws on the coordinate cubic tori of side
length $n=R/\eta$ in lattice units, with periodic gauge fields and fermions
antiperiodic in all three coordinate directions, and with
\begin{equation}\label{eq:cutoff-scales}
 \eta=L^{-N},\quad R=8hL^T,\quad
 h=L^{\lfloor\log_L\bar m^{-1}\rfloor},\quad
 q=e\eta^{1/2},\quad \boldsymbol\mu_N=\bar m\eta+m_{0,N}(e)
\end{equation}
with $N,T\in\Z_{\ge0}$ have volume-first limit points.
Every such limit point is carried by closed tempered two-forms,
satisfies Gaussian domination with constant one, and has tensor
Schwinger functions satisfying \textup{(E0$'$)}--\textup{(E3)} below.
Let $\mu$ be any such limit point and
$\mu=\int\mu_\xi\,\dd\nu(\xi)$ its decomposition into
translation-ergodic laws. For $\nu$-almost every $\xi$, $\mu_\xi$ is
carried by closed tempered two-forms and satisfies
\begin{equation}\label{eq:main-domination}
 \int e^{F(g)}\,\dd\mu_\xi(F)\le e^{\|g\|_{L^2}^2/2}
 \qquad(g\in\Sspace(\R^3;\Lambda^2)\text{ real}).
\end{equation}
Its Schwinger functions
$S_n^\xi(g_1\otimes\cdots\otimes g_n)
=\int\prod_{a=1}^nF(g_a)\,\dd\mu_\xi(F)$
extend to the complex tensor Schwartz spaces
$\Sspace_n=\Sspace(\R^{3n};(\Lambda^2\C^3)^{\otimes n})$
and obey the following conditions, with $S_0^\xi=1$.
\begin{enumerate}
\item[\textup{(E0$'$)}] For $n\ge1$ and $\Psi\in\Sspace_n$,
\begin{equation}\label{eq:main-growth}
 |S_n^\xi(\Psi)|\le2\,66^n\sqrt{n!}\,q_{6n}(\Psi),\qquad
 q_m(\Psi)=\max_{|\alpha|\le m}\sup_y
 (1+|y|^2)^{m/2}|\partial^\alpha\Psi(y)|,
\end{equation}
where the maximum is also over tensor components.
\item[\textup{(E1)}] For every Euclidean motion $g(x)=Qx+a$,
$Q\in O(3)$, $S_n^\xi(g\cdot\Psi)=S_n^\xi(\Psi)$, where
\[
 (g\cdot\Psi)(x_1,\ldots,x_n)
 =(\Lambda^2Q)^{\otimes n}\Psi(g^{-1}x_1,\ldots,g^{-1}x_n).
\]
\item[\textup{(E2)}] Let $r$ be reflection in a plane through the origin
and $H$ either closed half-space bounded by that plane.
If $\Psi_0\in\C$ and $\Psi_k\in\Sspace_k$, $1\le k\le M$,
have $\supp\Psi_k\subset H^k$, then
\begin{equation}\label{eq:main-rp}
 \sum_{k,l=0}^M S_{k+l}^\xi
 \bigl(\Theta_k\Psi_k\otimes\Psi_l\bigr)\ge0,
 \qquad \Theta_k\Psi_k=\overline{r\cdot\Psi_k},\quad
 \Theta_0c=\overline c.
\end{equation}
\item[\textup{(E3)}] $S_n^\xi$ is unchanged by simultaneous permutation
of the arguments and their tensor indices.
\item[\textup{(E4)}] For $\Psi\in\Sspace_k$ and $\Phi\in\Sspace_l$,
\begin{equation}\label{eq:main-cluster}
 S_{k+l}^\xi(\Psi\otimes\tau_a^{\otimes l}\Phi)
 \longrightarrow S_k^\xi(\Psi)S_l^\xi(\Phi)
 \quad\text{as }|a|\longrightarrow\infty,
\end{equation}
where
\[
 (\tau_a^{\otimes l}\Phi)(x_1,\ldots,x_l)
 =\Phi(x_1-a,\ldots,x_l-a).
\]
\end{enumerate}
Moreover, for a constant $C_f$ depending only on the fixed
construction, reference mass and probe, and independent of
$0<e\le e_\beta$, the limit point and $\xi$,
\begin{equation}\label{eq:main-cumulant}
 \bigl|\kappa_4^{\mu_\xi}(f)-e^4b(f)\bigr|
 \le C_f e^{4+\sigma},\qquad
 \kappa_4^{\mu_\xi}(f)\ge\tfrac12e^4b(f)>0.
\end{equation}
In particular, $\mu_\xi$ is non-Gaussian.
The full-measure component set is asserted for each fixed parent
law, charge and probe.
\end{theorem}

Here $\kappa_4^\lambda(f)$ is the fourth cumulant of $F(f)$ under a
law $\lambda$. For the Gaussian-dominated laws considered here,
$F(f)$ has mean zero, so
$\kappa_4^\lambda(f)=S_4^\lambda(f^{\otimes4})
-3S_2^\lambda(f^{\otimes2})^2$.
The charge threshold may depend on the fixed probe.
The factorization in \textup{(E4)} is a conclusion for the
translation-ergodic components.

The tensor action in \textup{(E1)} is essential to
\textup{(E2)}. For reflection in $x_1=0$, the components $F_{12}$ and
$F_{13}$ change sign, while $F_{23}$ does not. Replacing this action
by scalar reflection changes the positivity assertion.
The growth condition is stated on all tensor tests, including tests
meeting the diagonals; bounds on products of separately smeared fields
alone do not state \textup{(E0$'$)}.


\subsection{The proof and its relation to earlier work}

\Cref{thm:main} states the tensor form of the axioms of
Osterwalder and Schrader \cite{OS1973,OS1975}. Uniform Gaussian
estimates give tightness of the cutoff laws and convergence of their
moments along weakly convergent sequences. The lattice reflection
argument supplies positivity at the
three coordinate and six diagonal mirror directions; the comparison
with mirror-adapted tori requires taking the volume limit first.
These properties pass to every prescribed volume-first limit by
\cite{ChamberLattice,ContinuumLimits}. Reflection-positive lattice
gauge theories and their transfer matrices were developed by
Osterwalder and Seiler \cite{OsterwalderSeiler1978}; the chamber
crossings and the comparison of the two torus geometries are the
construction-specific inputs used here.

The nine mirror directions give the spectral inequalities
$H_j\ge|\mathsf P_i|$ for $i\ne j$. Together with Gaussian
domination, they bound the angular frequencies of the Schwinger
functions. Closedness gives the required tensor weight at two
points; reflection positivity and the regular-polygon argument then
give that weight at every order. The proof in \cref{sec:rotations}
also treats coincidences and excludes contact terms. We apply it
to almost every ergodic component, after obtaining the Gaussian
bound and the nine reflection-positive forms on a common component
set from \cite{ErgodicComponents}. The same spectral inequalities,
together with the local correlation estimate, give clustering in
each component. The supported-tensor argument
in \cref{sec:tensor-axioms} extends these conclusions to the tests
in \textup{(E0$'$)}--\textup{(E4)}, including closed-half-space
supports in reflection positivity.

The renormalization estimates select the volume-independent mass
and control the fourth-order response. Dimock \cite{DimockUVS2022}
proves ultraviolet stability at weak coupling in a fixed finite
volume using Ba{\l}aban's renormalization group. The block analysis
used here adapts this framework to chamber fermions and two
reference tori. Decay estimates for the local inverses control the
polymer activities; extraction of the vacuum and mass terms leaves
a normalized interaction that contracts under block averaging.
The nonlinear step and the common tuning theorem give controlled
trajectories for one mass and vacuum assignment on both reference
tori \cite{NormalizedContraction,NonlinearStep,CoverTuning}.
Terminal pressure estimates then enter the quantitative response
bound \cite{TerminalPressure,ReferenceResponse}. This control of
the interacting flow supplies the non-Gaussianity estimate, while
the Gaussian bounds supply the compactness and moment bounds used
in the limit passage.

The transfer from the reference tori uses chessboard estimates,
whose general framework is due to Fr\"ohlich, Israel, Lieb, and
Simon \cite{FILS1978}. The volume-transfer theorem of
\cite{VolumeTransfer} controls products of translated Wick
powers as well as the fourth cumulant. For a prescribed parent law,
these product bounds control invariant conditional expectations
and yield the cumulant estimate in almost every ergodic component
of that same law. The free fourth coefficient converges to $b(f)$,
and the continuum calculation in \cite{PositiveProbe} provides a
compact probe with $b(f)>0$. Intersecting the resulting full-measure
component set with the structural one, and reducing the
probe-dependent charge interval, gives \eqref{eq:main-cumulant}
along with all five tensor axioms.
The proof of \cref{thm:main} is completed in
\cref{sec:main-theorem-closure}.
\section{The finite laws and their continuum observables}\label{sec:regularization}

We use the all-flux chamber model of \cite[Section~2]{ChamberLattice}.
The following definitions fix its normalization and its relation to the
physical scales in \eqref{eq:cutoff-scales}. All cochain norms and
pairings sum over positively oriented cells. Reversing the orientation
changes the sign of a cochain.

\subsection{The chamber matrix and the flux sum}
Let $O_{\rm cub}$ be the signed permutation matrices in three dimensions.
A chamber label is $(c,g)$, where $c\in\Z^3$ and $g\in O_{\rm cub}$.
There are $48$ labels at each photon site. Let $s_0$ reverse coordinate
three, let $s_1,s_2$ exchange coordinates $(2,3)$ and $(1,2)$, and let
$r_{e_1}$ reverse coordinate one. With the Pauli matrices $\sigma_j$, put
$\gamma_j=\operatorname{diag}(\sigma_j,-\sigma_j)$ and
\[
 \begin{gathered}
 (w_0,w_1,w_2,w_3)=(6,6\sqrt2,6\sqrt2,6),\\
 (\mathbf n_0,\mathbf n_1,\mathbf n_2,\mathbf n_3)
 =(e_3,(e_2-e_3)/\sqrt2,(e_1-e_2)/\sqrt2,-e_1),\qquad
 M_i(g)=\tfrac12w_i(\rho+\gamma\cdot g\mathbf n_i).
 \end{gathered}
\]
For a bond phase $B$, the unit-mesh massless matrix is
\begin{equation}\label{eq:reg-dirac}
 \begin{split}
 (D^{\rm ch}[B]\psi)(c,g)={}&
 \sum_{i=0}^2M_i(g)[\psi(c,g)-\psi(c,gs_i)]\\
 &+M_3(g)[\psi(c,g)-e^{iB(c,c+ge_1)}
                         \psi(c+ge_1,gr_{e_1})].
 \end{split}
\end{equation}
On $C(n)=\Z^3/n\Z^3$ the bond fields are periodic and spinors
satisfy $\psi(c+ne_j,g)=-\psi(c,g)$. The massive matrix
$D^{\rm ch}[B]+\boldsymbol\mu$ is the matrix in
\cite[Section~2]{ChamberLattice} at diagonal parameter
$\boldsymbol\mu+\frac\rho2\sum_iw_i$, where $\boldsymbol\mu$ is
the total lattice mass. The chamber inner product has weight $1/48$;
this constant weight does not change a finite-dimensional determinant.

For $q>0$, put $r_q=2\pi/q$. A photon configuration is a pair
$(A,z)$ with $A$ a real one-cochain and $z$ a closed integer
two-cochain, modulo
\[
 (A,z)\sim(A+r_q k,z-dk),\qquad k\text{ an integer one-cochain}.
\]
Its curvature is $u=dA+r_qz$. With the representative
$A(b)\in[0,r_q)$, define
\begin{equation}\label{eq:reg-functional}
 Z_{C(n)}[G]=
 \sum_{\substack{z\in C^2(C(n);\Z)\\dz=0}}
 \int_{[0,r_q)^{3n^3}}
 G(A,z)e^{-\|u\|^2/2}
 \det(D^{\rm ch}[qA]+\boldsymbol\mu)\,\dd A.
\end{equation}
The determinant is nonnegative and $0<Z_{C(n)}[1]<\infty$ for
every real $\boldsymbol\mu$, by
\cite[Propositions~2.2--2.3]{ChamberLattice}.
We write $\E_n^{\rm af}=Z_{C(n)}[\,\cdot\,]/Z_{C(n)}[1]$.
Rectangular tori carry the same periodic photon and antiperiodic
fermion conventions. No flux class is excluded from this law.

At the physical cutoff, $q=e\sqrt\eta$ and
$\boldsymbol\mu=\boldsymbol\mu_N=\bar m\eta+m_{0,N}(e)$.
The physical field $F^\sharp$ is the centered step interpolation
of $\eta^{-3/2}u$, defined precisely in \cref{sec:continuum-limits}.
For a fixed mass sequence, $P_{N,T}$ denotes its law on $H_{-4}$.
Let $\mathcal V_N$ be the set of its weak limits as $T\to\infty$.
A volume-first limit is a weak limit of
$\mu_i\in\mathcal V_{N_i}$ with $N_i\to\infty$.

\subsection{Gaussian bounds and reflection quotients}\label{sec:probability}
For real Schwartz two-forms, Gaussian domination means
$\E_\lambda e^{F(f)}\le e^{\|f\|_2^2/2}$.
The estimates of \cite[Proposition~2.1]{ContinuumLimits} give
\[
 \E_\lambda F(f)=0,\qquad
 \E_\lambda|F(f)|^p\le c_p\|f\|_2^p,\qquad
 c_p=2\,2^{p/2}\Gamma(1+p/2).
\]
For complex tests the product bound has the additional factor
$2^{n/2}$ at degree $n$. Applied to differences, these estimates
extend $F$ from Schwartz tests to $L^2$, consistently in every
$L^p(\lambda)$, $p\ge1$. Expanding one factor at a time likewise
extends each joint moment as a bounded multilinear form on complex
$L^2$. These completions will be used for the plate insertions.
The same proposition gives entire finite-dimensional generating
functions, so their moments determine the finite-dimensional laws.

Write $h_\kappa=h_k e_{ij}$ for the real Hermite two-forms,
$k\in\Z_{\ge0}^3$, $i<j$, with eigenvalues
$\lambda_\kappa=2|k|_1+4$ for $1+|y|^2-\Delta$. Set
\[
 \|F\|_{-s}^2=\sum_\kappa\lambda_\kappa^{-2s}|F(h_\kappa)|^2,
 \qquad H_{-s}=\{F\in\mathcal S':\|F\|_{-s}<\infty\}.
\]
All weak convergence of laws on $H_{-4}$ uses its norm topology.
The Hermite compactness argument of \cite[Section~4]{ContinuumLimits}
shows that a closed bounded $H_{-3}$ ball is compact in $H_{-4}$.
The tensor estimate \cite[Proposition~6.1]{ContinuumLimits} gives a
continuous map $Q_\Psi\in L^2(\lambda)$ for $\Psi\in\mathcal S_n$,
equal to $\prod_aF(f_a)$ on elementary tensors, and the bound
\eqref{eq:main-growth} for its expectation.

For open $O$, let $\Sigma_O$ be generated by $F(f)$ with real
$f\in C_c^\infty(O;\Lambda^2)$. For a Euclidean motion $g$,
write $\widehat g A(F)=A(g^{-1}\cdot F)$ and $U(a)=\widehat{\tau_a}$.
Let $r_j$ reflect in $x_j=0$ and put
$\mathcal R_j=\widehat r_j$, $\Theta_j A=\overline{\mathcal R_jA}$.
For a translation-invariant law that is invariant and reflection
positive at this plane, the positive form
\[
 b_j(A,B)=\E_\lambda[(\Theta_jA)B],\qquad
 A,B\in L^2(\Sigma_{\{x_j>0\}},\lambda),
\]
has quotient completion $\mathcal H_j$, quotient map $[A]_j$,
and vector $\Omega_j=[1]_j$. The reconstruction and spectral cone
are \cite[Propositions~2.1 and~3.1]{ErgodicComponents}:
\[
 e^{-tH_j}V_j(y)[A]_j=[U(te_j+y)A]_j,\qquad
 t\ge0,\quad y\perp e_j,\qquad H_j\ge|\mathsf P_i|\quad(i\ne j).
\]
The cone conclusion uses invariance and reflection positivity at
$x_j+x_i=0$. These are the precise reconstruction inputs used in
\cref{sec:rotations}.

\subsection{Block scales and the selected pair}\label{subsec:integration-block-data}
The block construction uses an odd $L$, an integer power $M$ of $L$,
and the primitive and activity conventions of
\cite{MaxwellBlock,FermionBlock,NormalizedContraction,NonlinearStep,CoverTuning}.
Their constants will be chosen in \cref{sec:ordered-constants}.
For the reference mass in \cref{thm:main}, write
\[
 K=N+\log_Lh,\qquad a_j=L^{j-N},\quad e_j=e\sqrt{a_j},
 \quad\mu_j=\bar m a_j,\quad v_j=-\log e_j.
\]
Let $r_j$ be the least nonnegative integer power of $L$ at least
$v_j^2$. Put
\[
 R_{\rm top}=hMr_K,\qquad R^*=R_{\rm top}/L,\qquad
 \mathcal T(e)=\{hL^t:t\in\Z_{\ge0},\ hL^t\le R_{\rm top}\}.
\]
For $R\in\mathcal T(e)$, the stop is the first level
$k_*(R)$ with $\ell_j=Mr_j$, where $\ell_j=R/a_j$.
The finite threshold $N_{\min}(R,e)$ and the existence and
properties of the stop are given in \cite[Proposition~6.1]{CoverTuning}.
A star on a running quantity means evaluation at this stop.
Write $N_{\rm ref}(e)=N_{\min}(R^*,e)$ and $n_*=R^*/\eta$.

Under the numerical conditions verified in \cref{prop:ordered-constants}, the common-pair theorem
of \cite[Theorem~A]{CoverTuning} selects
$P_N(e)=(m_{0,N}(e),\varepsilon^v_{0,N}(e))$ on the top torus,
with zero terminal mass shift and vacuum term there. The same pair
is used on the smaller reference tori, for every twist and source.
Its scalar bounds and controlled trajectories are conclusions of
that theorem. Below the top definition threshold, only
$m_{0,N}(e)$ is set to zero. No trajectory is used at those cutoffs.
The construction of this pair requires no choice of the probe.

The early zero-flux source identities in
\cite[Section~3]{ReferenceResponse} define the matter normalization
used by \cite{TerminalPressure}. The terminal measure and pressure
estimates of the latter then enter the quantitative response theorem
of \cite{ReferenceResponse}. In \cref{sec:ordered-constants} we
first verify the construction and terminal scalar conditions.
We then apply the response theorem and the volume-transfer estimates
to choose the probe-dependent interval.
\section{Volume and continuum limits}\label{sec:continuum-limits}

Gaussian domination gives compactness of the curvature laws at all
cutoffs. Diagonal reflection positivity uses a further identification:
at a fixed spacing, a volume limit is the step-field image of an
infinite lattice law. This permits the nine-mirror theorem to be
applied before the spacing is removed.

\subsection{The test map and its bounds}

Write $\eta=L^{-N}$, $R=8hL^T$, and $n=R/\eta$, as in
\eqref{eq:cutoff-scales}. A plaquette of the unit covering lattice is
$\sigma=(x;\mu\nu)$, with $\mu<\nu$ and center
$y_\sigma=x+(e_\mu+e_\nu)/2$. Its averaging cube is
\[
 Q_\eta(\sigma)=\eta\bigl(y_\sigma+[-1/2,1/2)^3\bigr).
\]
For each fixed orientation these cubes partition $\R^3$.
If $u$ is a real lattice curvature, its physical step field is
$\eta^{-3/2}u(\sigma)$ on $Q_\eta(\sigma)$ in component $\mu\nu$.
For a Schwartz two-form $f$, define
\begin{align}
 f_{[\eta]}(\sigma)
   &=\eta^{-3/2}\int_{Q_\eta(\sigma)}f_{\mu\nu}(y)\,\dd y,
                                                   \label{eq:limit-test-cover}\\
 f^R(y)&=\sum_{k\in\Z^3}f(y+Rk),\notag\\
 f_\eta([\sigma])
   &=\eta^{-3}\int_{Q_\eta(\sigma)}f^R_{\mu\nu}(y)\,\dd y.
                                                   \label{eq:limit-test-torus}
\end{align}
Here $[\sigma]$ is the plaquette class modulo $n\Z^3$.
The torus cochain norm is $|v|_\eta^2=\eta^3\sum_{[\sigma]}|v([\sigma])|^2$.
The covering cochain norm $\|\cdot\|$ is the unweighted $\ell^2$ norm.
For the periodic extension $\widetilde u$ of $u$,
\begin{equation}\label{eq:limit-step-pairing}
 F^\sharp(f)=\eta^{3/2}\sum_{[\sigma]}f_\eta([\sigma])u(\sigma)
            =\sum_\sigma f_{[\eta]}(\sigma)\widetilde u(\sigma).
\end{equation}
The normalization in the first equality combines the cell volume
$\eta^3$ with the physical field factor $\eta^{-3/2}$.

The estimates needed below are
\begin{align}
 \|f_{[\eta]}\|&\le\|f\|_2, \label{eq:limit-cover-norm}\\
 \sum_\sigma|f_{[\eta]}(\sigma)|(1+|x(\sigma)|_\infty)
 &\le\eta^{-3/2}\int_{\R^3}\sum_{\mu<\nu}|f_{\mu\nu}(y)|
                         (2+|y|_\infty/\eta)\,\dd y, \label{eq:limit-cover-sum}\\
 |f_\eta|_\eta&\le\|f^R\|_{L^2(\R^3/R\Z^3)}\le p_R(f),\label{eq:limit-torus-norm}
\end{align}
where, with $Q_R=[-R/2,R/2)^3$,
\begin{equation}\label{eq:limit-periodic-seminorm}
 p_R(f)=\|f\|_2+
        26R^{-2}\||y|^2f\|_{L^2(\R^3\setminus Q_R)}.
\end{equation}
One also has $|f_\eta|_\eta\le\|f\|_2$ if the support of $f$ lies
in a translate of an open cube of side $R$. For $R\ge1$,
$p_R(f)\le p_1(f)$, and $p_R(f)\to\|f\|_2$ as $R\to\infty$.
The seminorm $p_1$ is continuous on Schwartz space.

The first bound follows from Cauchy--Schwarz on each averaging
cube. On that cube $1+|x(\sigma)|_\infty\le2+|y|_\infty/\eta$,
which gives the second. Absolute convergence permits unfolding:
\[
 \eta^{3/2}f_\eta([\sigma])
       =\sum_{k\in\Z^3}f_{[\eta]}(\sigma+nk).
\]
This proves \eqref{eq:limit-step-pairing}. Applying the cube
Cauchy--Schwarz inequality to $f^R$ gives the first inequality in
\eqref{eq:limit-torus-norm}. If the support lies in one open
$R$-cube, its distinct $R\Z^3$ translates are disjoint, so
$\|f^R\|_{L^2(\R^3/R\Z^3)}=\|f\|_2$.

For the remaining estimate put $a_k=\|f\|_{L^2(Q_R+Rk)}$.
Minkowski's inequality bounds the periodized norm by $\sum_k a_k$.
On $Q_R+Rk$, $k\ne0$, one has $|y|\ge R|k|_\infty/2$; hence
\[
 \sum_{k\ne0}a_k
 \le4R^{-2}\left(\sum_{k\ne0}|k|_\infty^{-4}\right)^{1/2}
                   \||y|^2f\|_{L^2(\R^3\setminus Q_R)}.
\]
The number of points with $|k|_\infty=j$ is $24j^2+2$, so
the sum in parentheses is $4\pi^2+\pi^4/45<41.65$ and its
square root times four is less than $26$.
This proves \eqref{eq:limit-torus-norm} and
\eqref{eq:limit-periodic-seminorm}. The comparison with $p_1$
and the limit of $p_R$ follow by inclusion of the cubes and
integrability of $|y|^4|f(y)|^2$.

For the Hermite two-forms of \cref{sec:probability},
\begin{equation}\label{eq:limit-Hermite-norm}
 p_1(h_\kappa)\le31\lambda_\kappa.
\end{equation}
Here is the weighted estimate behind this bound. Let
$\mathsf M=1+|y|^2$, $P_j=-i\partial_j$, and
$\mathsf A=\mathsf M-\Delta$. Since
$[P_j,[P_j,\mathsf M]]=-2$,
\[
 \|\mathsf A v\|_2^2
 =\|\mathsf Mv\|_2^2+\|\Delta v\|_2^2
      +2\sum_j\langle P_jv,\mathsf MP_jv\rangle-6\|v\|_2^2
 \ge\|\mathsf Mv\|_2^2-6\|v\|_2^2.
\]
The Hermite eigenvalues satisfy $\lambda_\kappa\ge4$, and therefore
$\|\mathsf Mv\|_2^2\le(11/8)\|\mathsf Av\|_2^2$.
Taking $v=h_\kappa$ gives
$p_1(h_\kappa)\le1+26\sqrt{11/8}\lambda_\kappa
\le31\lambda_\kappa$.

The test map respects the affine cubic symmetries. For a signed
permutation and lattice translation $g$ of the unit lattice, let
$\mathfrak g(y)=\eta g(y/\eta)$. With the exterior-square action
on two-forms and the cochain pull-back,
\begin{equation}\label{eq:limit-test-equivariance}
 (\mathfrak g\cdot f)_{[\eta]}=(g^{-1})^*f_{[\eta]},
 \qquad
 (\mathfrak g\cdot f)_\eta=(g^{-1})^*f_\eta.
\end{equation}
Indeed, $g$ sends centers to centers and the averaging cubes to
averaging cubes up to null boundaries. Its orientation sign is
the sign of $\wedge^2\vec g$ on the corresponding plaquette.
Changing variables gives the first formula; periodization and
the bijection $k\mapsto\vec gk$ give the second.

\subsection{Compactness and convergence of moments}

Let $P_{N,T}$ be the law of $F^\sharp$ for the all-flux chamber
measure. The mass assignment may be any real sequence independent of
$T$. Proposition~\extref{5.1}{ChamberLattice} and
\eqref{eq:limit-step-pairing} give
\begin{equation}\label{eq:limit-cutoff-G}
 \E_{N,T}e^{F^\sharp(f)}
       \le e^{|f_\eta|_\eta^2/2}\le e^{p_R(f)^2/2}
       \qquad(f\text{ real}).
\end{equation}
The random-linear-form proof of Proposition~\extref{2.1}{ContinuumLimits}
applies here with the Hilbert seminorm $f\mapsto|f_\eta|_\eta$.
In particular its variance bound and \eqref{eq:limit-Hermite-norm}
give
\begin{equation}\label{eq:limit-uniform-tightness}
 \E_{N,T}\|F^\sharp\|_{-3}^2
 \le\sum_\kappa\lambda_\kappa^{-6}p_1(h_\kappa)^2
 \le961\sum_\kappa\lambda_\kappa^{-4}<61.
\end{equation}
Tonelli's theorem justifies the first sum. Every finite lattice
cochain has a step field in $H_{-3}$: Cauchy--Schwarz in
\eqref{eq:limit-step-pairing} bounds its $\kappa$-th Hermite
coordinate by $31\lambda_\kappa\|u\|_{\ell^2(C(n)^{(2)})}$,
with the unweighted norm on the unit-lattice torus plaquettes.
The step map from this finite-dimensional cochain space to $H_{-3}$
is consequently continuous.

The compact embedding recalled in \cref{sec:probability}
now applies. The laws $P_{N,T}$ are uniformly tight on $H_{-4}$,
and their fixed-spacing volume limits are uniformly tight there.
Thus the sets $\mathcal V_N$ and the set of volume-first limits
are nonempty.

For a joint sequence $N_k,T_k\to\infty$ suppose
$P_{N_k,T_k}\to\mu$ weakly on $H_{-4}$. Truncating $e^{F(f)}$
above by a constant, passing to the weak limit, and then removing
the truncation gives
\begin{equation}\label{eq:limit-continuum-G}
 \E_\mu e^{F(f)}\le e^{\|f\|_2^2/2}.
\end{equation}
The limiting constant follows from $p_{R_k}(f)\to\|f\|_2$.
For a fixed finite family of complex tests $f_1,\ldots,f_s$, define
\[
 Z_k(z)=\E_{N_k,T_k}\exp\left(\sum_a z_aF^\sharp(f_a)\right),
 \qquad
 Z(z)=\E_\mu\exp\left(\sum_a z_aF(f_a)\right).
\]
The evaluations $F\mapsto(F(f_a))_a$ are continuous on $H_{-4}$.
Moreover,
\[
 \E_{N_k,T_k}\left|\exp\left(\sum_a z_aF^\sharp(f_a)\right)\right|^2
 \le\exp\left(2p_1\left(\operatorname{Re}\sum_a z_af_a\right)^2\right).
\]
This bound is uniform in $k$ and on compact sets of $z$.
Uniform integrability therefore gives pointwise convergence
$Z_k\to Z$. The same exponential bounds give local uniform
bounds on the entire functions $Z_k$. Cauchy's estimates make them
equi-Lipschitz on compact sets, so their convergence is locally
uniform. Cauchy's formula then passes every derivative to the limit.
In particular all joint moments converge.

On each compact subset of
$\{\sum_a|z_a|\|f_a\|_2<2/5\}$, the test-map bounds give
$\sum_a|z_a||(f_a)_{\eta_k}|_{\eta_k}<2/5$ for all large $k$.
The zero-free estimate \exteq{17}{ContinuumLimits}, with this
seminorm, bounds $|Z_k-1|$ by $1/2$. The principal logarithms
therefore converge locally uniformly near zero, and their
derivatives give convergence of every joint cumulant.

\subsection{The prescribed parent and its components}

The finite-periodic moment passage just proved applies to a
diagonal realization of any prescribed volume-first limit.
The fixed-spacing recovery in \cite[Proposition~7.3]{ContinuumLimits}
identifies each intermediate volume law with the interpolation
of an infinite-lattice law. Its nine reflection forms then pass
to the continuum limit by \cite[Corollary~7.4]{ContinuumLimits}.
We record the resulting input to the construction.

\begin{proposition}\label{prop:limit-nine-mirrors}
Let $\mu$ be any volume-first limit for a real, volume-independent
mass assignment. Then $\mu$ is carried by closed tempered two-forms,
is translation invariant and cubic invariant, and satisfies
\eqref{eq:limit-continuum-G}. It is invariant and reflection
positive at each plane in the nine families
\[
 x_j=c\quad(1\le j\le3),\qquad
 x_j\pm x_k=c\quad(1\le j<k\le3),\qquad c\in\R.
\]
There is a sequence of finite laws $P_{N_i,T_i}\to\mu$ with
$N_i,T_i\to\infty$ along which every fixed finite-family moment
and cumulant converges. In the translation-ergodic decomposition
of $\mu$, a single full-measure set of components has the same
closedness, Gaussian bound, cubic symmetries and nine reflection
forms. Its component laws have clustering for all $L^2$ observables
in all directions.
\end{proposition}

\begin{proof}
The parent-law conclusions and the weak finite-law diagonal are
\cite[Corollary~7.4]{ContinuumLimits}; its hypotheses hold by
the all-flux lattice theorem of \cite[Theorem~A]{ChamberLattice}
and \eqref{eq:limit-cutoff-G}. The preceding moment passage
applies to that diagonal. The component conclusions are
\cite[Theorem~A]{ErgodicComponents}, with closedness taken as the
specified event of full parent measure. That theorem gives all
real Schwartz Gaussian bounds simultaneously and the full
translation action on the same component set.
\end{proof}

\section{Choice of the construction constants}\label{sec:ordered-constants}

The estimates impose conditions on fixed constants and on the
running charge. The constants are chosen before the charge, mesh,
volume and twist. We first choose the polymer decay and the two
geometric scales; the charge bounds are then imposed for every
smaller charge. The local inverse constants are fixed before $M$,
whereas the weakened kernel constants may depend on $M$.
Neither group depends on the running charge. We verify this order
before imposing the small-charge conditions. The final subsection
uses the resulting construction to make the additional probe-dependent
restriction needed for the fourth cumulant.

\begin{proposition}\label{prop:ordered-constants}
Fix $0<\rho<1$ and the standard chamber weights.
The constants of the photon and fermion primitives,
small-field step, cover comparison and tuned flow can
be chosen in the order described below, with positive
charge ceilings $\bar e_0,\bar e_1$ independent of the mesh,
volume, twist and reference mass.
For every $0<\bar m\le1$, the physical ceiling $e_G$
in \eqref{eq:constants-physical-charge} is positive.
The terminal conditions \exteq{43}{TerminalPressure}
hold on the top and reference tori throughout
$0<e\le e_G$.
\end{proposition}

\subsection{The geometric scales and the positive exponent}

\begin{proof} Set $\bar\varepsilon=1/200$, $c_X=54$,
$\kappa_0=54\log72$, $K_{\mathrm{pol}}=72^{54}/36$,
$c_{\mathrm{rb}}=54K_{\mathrm{pol}}$ and $c_{\mathrm{fl}}=61c_XK_{\mathrm{pol}}$.
Here $K_{\mathrm{pol}}$ is the scalar polymer-sum constant of
\cite[Appendix~B, (81)]{NormalizedContraction}. Choose
\[
 1-\bar\varepsilon<\alpha<1,\qquad
 \kappa\ge7\kappa_0+7.
\]
An odd integer $L\ge5$ is required to satisfy
\begin{align}
 L^{1/8-3\bar\varepsilon/2}&\ge2400,\qquad
 c_{\mathrm{fl}}L^{-3/8+6\bar\varepsilon}\le1/16,\notag\\
 L(\kappa-6\kappa_0-4)-1&\ge\kappa,\notag\\
 c_{\mathrm{rb}}L^3
       \exp\bigl(-L(\kappa-\kappa_0-1)+\kappa\bigr)
                                                   &\le1/2.
                                      \label{eq:constants-L}
\end{align}
These are simultaneous lower-bound conditions on $L$.
The powers $1/8-3\bar\varepsilon/2$ and
$3/8-6\bar\varepsilon$ are positive, as are
$\kappa-6\kappa_0-4$ and $\kappa-\kappa_0-1$.
The last expression tends to zero as $L\to\infty$.
Thus every sufficiently large odd $L$ satisfies
\eqref{eq:constants-L}.

After $L$ is fixed, the primitive construction fixes
$a'$, $C_V$, $\gamma_w$ and $\gamma_P$, with $a',C_V>0$
and $\gamma_w,\gamma_P>0$. Put
$N_c=(2L+1)^3$ and
$\widehat\kappa=N_c\kappa+\log(6c_X)>0$.
The photon and fermion
constants have no dependence on $M$ at this stage.
Put $\kappa_w=\kappa+\kappa_0+3$, $\alpha_0=1/54$ and $\kappa'=\kappa_w+\kappa_0$ and
\[
 M_F=\max\{4L+1,\ (108\pi C_Va')^2,\
                         108\pi C_Va'\exp(27\kappa')\}.
\]
Choose $M$ to be a power of $L$ satisfying
\begin{equation}\label{eq:constants-M}
 \begin{split}
 M&\ge L^2,\qquad M>4L,\qquad \gamma_wM\ge16,\\
 M^{\alpha_0}&\ge1+\exp(\widehat\kappa),\qquad
 M\ge M_F,\qquad 804\alpha_0\log M\le\gamma_PM.
 \end{split}
\end{equation}
All these conditions hold for every sufficiently large
power of $L$. For the last one, $\log x\le\sqrt x$
on $x\ge1$ gives the sufficient lower bound
$M\ge(804\alpha_0/\gamma_P)^2$.
The definition of $M_F$ gives the two fermion
conditions in \exteq{46}{FermionBlock},
because $27\alpha_0=1/2$.
The last inequality gives \exteq{38}{MaxwellBlock}.
The first four give the geometric separation, decay
and parameter-radius conditions used in localization.

The exponents can be fixed independently of the charge.
Take $r=2$, $p_0=9$, $p=19$ and
$1-\alpha<\varepsilon\le\bar\varepsilon$, and take $s=1/2$.
Then $p_0>3r+2$ and $p>2p_0$, while
\[
 \mathfrak e=s(1/4-9\varepsilon)-12\varepsilon
                      \ge17/400.
\]
For $N_*=8L^3$ and $a_L=(L^3-1)/(L-1)^3$,
\[
 \frac{2a_L}{N_*}
     <\frac1{4(L-1)^3}\le\frac1{256}<\frac{17}{400}.
\]
It follows that $\sigma=\mathfrak e-2a_L/N_*>0$.
In particular the positive exponent is fixed before
the charge is chosen.

\subsection{Calibration of the primitive and terminal bounds}

The chamber coercivity, heavy-sector gap and local inverse
construction fix $a'$ and $\gamma^F>0$ before $M$ is chosen;
see \exteq{34}{FermionBlock} and
\exteq{43}{FermionBlock}. The photon construction
likewise fixes $C_P^{\mathrm{ph}}$ and $\gamma_P>0$ before $M$;
its weakened families are considered only after
\exteq{38}{MaxwellBlock} holds.
The sums $c(\gamma)=\sum_{z\in\Z^3}
\exp(-\gamma|z|_\infty)$ are finite for $\gamma>0$.
Thus the constants used to choose $M$ in
\eqref{eq:constants-M} have already been fixed.

Keep $b=1$, $b_-=1-L^{-1}$ and $b_+=1$ from the inverse
construction; these are already fixed before $M$.
The coarse resolvent and its auxiliary block matrix give
\[
 C_{\mathrm{By}}=(1+2\exp(1))\sqrt2\max(1,b_-^{-1}),
 \qquad
 C_\Gamma=\frac{2(1+2\exp(1))}{b_-}
                  +128\exp(1)(1+2\exp(1))^2a'.
\]
The latter is the constant in
\exteq{62}{FermionBlock}. Define
\begin{equation}\label{eq:constants-fermion-calibration}
 \begin{split}
 C_1=\max\{&
 304\exp(1)\max(1,b)a'M^{27\alpha_0},
 C_\Gamma M^{27\alpha_0},\\
 &32\exp(1)\max(1,b^2)a'\exp(27\kappa')\},\\
 C^F=2\max\{&
 C_1,C_{\mathrm{By}},304\exp(1)\max(1,b)a',
 C_\Gamma,\exp(1)b^2(32a')^2\}.
 \end{split}
\end{equation}
The factor two converts Euclidean spinor bounds to
component bounds. The walk and coarse-family estimates
\exteq{47}{FermionBlock}--\exteq{48}{FermionBlock}
and \exteq{61}{FermionBlock}--\exteq{62}{FermionBlock}
are bounded by $C^F$ in this convention. The last entry
also bounds the compressed quadratic kernel.
At level zero the normalized kernel has the same bound,
and $M^{8\alpha_0}\le M^{27\alpha_0}$ bounds its diagonal
field multiplier.
The unweakened all-pair field bound
\exteq{60}{TerminalPressure} is at most
$608\exp(1)\max(1,b)a'\le C^F$; at zero its constant is
at most four, also below $C^F$.

Take the common primitive constants
\[
 C^P=116C_P^{\mathrm{ph}}M^{402\alpha_0},\qquad
 C_w^{(1)}=\max\{C^F,C^P,4\exp(2\gamma_w)M^{27\alpha_0}\},
 \qquad C_w=64K_{\mathrm{pol}}C_w^{(1)}.
\]
These are precisely \exteq{23}{CoverTuning}.
The top systems and their projections therefore satisfy
the same bounds at every required level and side.
The determinant, photon-domain and localization constants
are given in Sections~\extref{6.2}{FermionBlock},
\extref{6}{MaxwellBlock} and
\extref{5.1}{NonlinearStep}. In particular,
\[
 \begin{aligned}
 c_C&=3L^{9/2},\qquad C_W=28C_w^2c_\gamma,\\
 C_d=16c_XC_w
       &\max\{\tfrac\pi2b\exp(1)(1+2\exp(1)),
                                      2+58\exp(1)+b\exp(1)\}.
 \end{aligned}
\]
The kernel-sum constant $c_\gamma$ is bounded by
$2L^3\exp((2L+1)\gamma_w)c(\gamma_w)$;
the local-gauge constant $C_B$ depends only on
$C_P^{\mathrm{ph}},\gamma_P,L,M$.
All these constants are finite after $M$ is fixed.

For the terminal norms fix
\begin{equation}\label{eq:constants-terminal-calibration}
 C_0=\max\{C_P^{\mathrm{ph}}c_{\mathrm{bd}}(\gamma_P),
                 C^F,4b_++16C^Fc(\gamma^F)\},
 \qquad c_{\mathrm{bd}}(\gamma)=3\exp(3\gamma)c(\gamma).
\end{equation}
The photon row sums, both field-map row norms and the
terminal Dirac form in
Section~\extref{6.2}{TerminalPressure} are then
bounded by this $C_0$. In particular $C_0\ge1$.
The polymer sum has $1\le c_\kappa\le K_{\mathrm{pol}}$.
The free Bloch construction fixes the twist constants
$C_\Sigma<\infty$ and $\gamma_\Sigma>0$ independently of
the running charge, mass in its prescribed range, mesh
and side; see \exteq{37}{TerminalPressure}.

The differentiated-response constants are now fixed by
the terminal inverse and field-map estimates:
$\mathcal S=C^F(24/\gamma^F)^3$,
$C_K=226\mathcal S^2$, $C_K'=1328\mathcal S^2$ and
$C_H=203M(C^F)^2$. The real covariance coefficients
$g,h_G,g_1,\alpha_G$ and its rate $\vartheta_0$ are
defined with \exteq{74}{TerminalPressure}.
The insertion and mixed-logarithm coefficients are
$C_V^{\mathrm{ins}}=10^3C_H^2(24/\gamma^F)^{12}K_{\mathrm{pol}}$ and
$C_W^{\mathrm{mix}}$ in \exteq{109}{ReferenceResponse}.
Their inverse-decay rates are positive and their
coefficients finite because $b_->0$,
$\mu_-=L^{-3}>0$ and $\gamma^F>0$.
They depend on $L,M$ and the fixed chamber data;
they have no dependence on the later charge or probe.
The probe-dependent constants will be chosen last.

\subsection{Uniform running-charge bounds}

All coefficients of the small-charge conditions are
fixed after the preceding choices. The fluctuation substitution
constant from \cite[Section~1]{NonlinearStep} is
\[
 C_T=2\exp(3/2)
       [1+b_-^{-1}(54L\exp(6)+12\exp(1)+1)].
\]
We write $\bar e_0,\bar e_1$ for the fixed charge ceilings,
to distinguish them from the level-zero running charge $e_0$.
They are the ceilings denoted $e_0,e_1$ in the terminal and
response estimates, and agree with the barred convention in
\cite{CoverTuning}. The elementary
estimate used to meet them simultaneously is as follows.
For finitely many $A_i,d_i,\lambda_i$ with
$A_i,d_i\ge0$ and $\lambda_i>0$, and for $b>0$, put
\begin{equation}\label{eq:constants-tail-threshold}
 V(A,d,\lambda;b)
   =\max_i\frac2{\lambda_i}
       \log^+\left[
         \frac{nA_i}{b}
         \left(\frac{2d_i}{\lambda_i\exp(1)}\right)^{d_i}
       \right].
\end{equation}
Here $n$ is the number of terms, $\log^+0=0$
and a factor with $d_i=0$ is interpreted as one.
For all $v\ge V(A,d,\lambda;b)$ with $v>0$,
\begin{equation}\label{eq:constants-tail-bound}
 \sum_i A_iv^{d_i}\exp(-\lambda_iv)\le b.
\end{equation}
Indeed $v^d\exp(-\lambda v/2)$ has maximum
$(2d/(\lambda\exp(1)))^d$ for $d>0$.
For $d=0$ use $\exp(-\lambda v)\le\exp(-\lambda v/2)$.
The threshold makes each resulting term at most $b/n$.

The primitive domains, determinant, adjustment and nonlinear
step require \exteq{43}{MaxwellBlock},
\exteq{8}{NormalizedContraction},
\exteq{21}{NormalizedContraction},
\exteq{1}{NormalizedContraction},
\exteq{81}{FermionBlock} and
\exteq{3}{NonlinearStep}.
Their charge conditions are obtained from the following
finite families for $v=-\log e_k$.
Write $T(A,d,\lambda)=Av^d\exp(-\lambda v)$.
The polynomial-exponential conditions are
\begin{align}
 T(1,p_0,\varepsilon/4)&\le1,&
 T(1,0,5\varepsilon/2)&\le\tfrac12,\notag\\
 T(c_{\mathrm{fl}},0,\varepsilon)&\le1,&
 T(3M,0,\varepsilon)&\le1,\notag\\
 T(2C_w^2c_\gamma^2c_C,0,\varepsilon/4)&\le1,&
 T(16c_Xc_\gamma C_wC_T^2M^3,0,\varepsilon/2)&\le1,\notag\\
 T(96M^3,0,1/4)&\le1,&
 T(32C_wC_T,0,3\varepsilon/2)&\le1,\notag\\
 T(88M^3,0,1/4-6\varepsilon)&\le1,&
 T(112C_w^2c_\gamma,0,2\varepsilon)&\le3,\notag\\
 T(c_S,0,1/4-5\varepsilon)&\le1,&
 T(c_SL^3,0,\varepsilon)&\le1,\notag\\
 T(36c_X^2K_{\mathrm{pol}}^2,0,1/4-5\varepsilon)&\le1,&
 T(1,0,\varepsilon)&\le\tfrac14,\notag\\
 T(3M,0,1/4+\varepsilon)+T(3M,0,2\varepsilon)&\le1,\notag\\
 T(1,0,1/4)+T(1,0,\varepsilon)&\le1,\notag\\
 T(160L^{1+3\varepsilon/2}M,0,\varepsilon)&\le1,&
 T(90c_XLM,0,\varepsilon)&\le1,\notag\\
 T(9ML^\varepsilon,0,2\varepsilon)&\le1,\notag\\
 T(2M,1,\varepsilon)+T(32M,0,\varepsilon)&\le1,\notag\\
 T(8C_B,p,3/4-3\varepsilon)&\le1,&
 T(8,0,3/4-7\varepsilon/2)&\le1,\notag\\
 T(6L,0,\varepsilon)&\le1,&
 T(4C_Va',0,\varepsilon)&\le1,\notag\\
 T(4,0,1/4+2\varepsilon)+T(2,0,3\varepsilon)&\le1,\notag\\
 T(8c_{\mathrm{rb}}L^3C_dM^3,0,2\varepsilon)&\le1.
                                      \label{eq:constants-running-inventory}
\end{align}
Here $c_S=10^5c_X^5K_{\mathrm{pol}}^5\exp(6\kappa)$.
The inequality with $96M^3$ controls the first step;
the one with $C_T^2M^3$ controls the positive-level
fluctuation substitution. The conditions with $c_S$
and $36c_X^2K_{\mathrm{pol}}^2$ control the connected logarithm and its scaling.
The domain conditions occur with target $1/4$,
with rates $(1/4,\varepsilon)$, with coefficients
$8C_B,8$, and with coefficients $4,2$.
The remaining displayed charge hypotheses above are
individual inequalities of this inventory;
\exteq{83}{NonlinearStep} follows from the
condition containing $c_S$.

Every decay rate in the inventory is positive for
$0<\varepsilon\le1/200$. The only additional
running-charge restrictions in this group are
\[
 v^{2p_0}/2-7v\ge\log12,\qquad v^2\ge10L,\qquad v\ge9.
\]
They hold respectively for $v\ge15$,
$v\ge\sqrt{10L}$ and $v\ge9$.
Taking the maximum of these thresholds and the finitely
many values in \eqref{eq:constants-tail-threshold}
gives a finite $V_0$. Choose a running-charge ceiling
$0<\bar e_0<\exp(-V_0)$ so that every such condition
holds for every $v\ge-\log \bar e_0$, equivalently for
every running charge in $(0,\bar e_0]$.
The set of permitted ceilings is downward closed.

The terminal-scale conditions require a second ceiling
$\bar e_1\le\min(\bar e_0,1/3)$. Write $\ell(v)=LMv^r$.
They follow from
\begin{align}
 (c_{\mathrm{fl}}+2)\exp(-\varepsilon v)
                      &\le\tfrac1{16}(1-L^{-\varepsilon/2}),\notag\\
 8C_0^2\ell(v)^6(b_-^{-1}+L\ell(v))
                      &\le\exp(\varepsilon v/2),\notag\\
 4(2C_0+c_\kappa+5c_{\mathrm{fl}}+6)\ell(v)^3
                      &\le\exp(\varepsilon v),\notag\\
 \exp(-(1/4-9\varepsilon)v)&\le\tfrac18,\notag\\
 4C_0\ell(v)^2v^p&<\exp(\varepsilon v),\notag\\
 12\ell(v)^3\left(\exp(-v^{p_0}/2)
                    +\exp(-v^{2p}/(2C_0^2))\right)
                      &\le\exp(-v),\notag\\
 C_\Sigma\exp(-\gamma_\Sigma Mv^r)&\le\exp(-v),\notag\\
 L^2\exp(-2v)\ell(v)&\le1,\notag\\
 6N_*Mv^r\exp(-v)(v\exp(-v)+\tfrac12)&\le\pi^2.
                                      \label{eq:constants-terminal-inventory}
\end{align}
These are the scalar hypotheses of
\exteq{48}{CoverTuning},
\exteq{43}{TerminalPressure} and
\eqref{eq:transfer-tail-condition}, after substituting
\exteq{51}{CoverTuning}.
Every polynomial-exponential rate here is positive for
$0<\varepsilon\le1/200$. The reference flux condition expands as
\[
          6N_*Mv^{r+1}\exp(-2v)+3N_*Mv^r\exp(-v)\le\pi^2.
\]
The superlinear tails reduce to these bounds:
for $v\ge\max(4,4C_0^2)$,
\[
 v-\tfrac12v^{p_0}\le-v,\qquad
 v-\frac{v^{2p}}{2C_0^2}\le-v.
\]
Hence
\[
 12\ell(v)^3
 \left[\exp\left(v-\tfrac12v^{p_0}\right)
       +\exp\left(v-\frac{v^{2p}}{2C_0^2}\right)\right]
                \le24(LM)^3v^{3r}\exp(-v).
\]
For the twist tail, if
$v\ge\max(1,2/(\gamma_\Sigma M),\log^+C_\Sigma)$,
then $C_\Sigma\exp(v-\gamma_\Sigma Mv^r)\le1$.
Put $\mathrm{lg}=\tfrac12\log L$.
The scale-ratio condition follows for $v\ge2r$:
\begin{equation}\label{eq:constants-scale-ratio}
 \left(1+\frac{3\mathrm{lg}}{v}\right)^r
        \le\exp\left(\frac{3r\mathrm{lg}}{v}\right)
        \le L^{3/4}<L.
\end{equation}
Together with $v\ge8$, $v\ge31L+17$ and
$p_0v^{p_0-1}\ge28$, these bounds give a finite
threshold $V_1$. Its maximum includes $2r$.
Let $\mathcal E_1$ consist of the numbers
$e'\in(0,\min(\bar e_0,1/3)]$ for which the geometric
conditions and \eqref{eq:constants-terminal-inventory}
hold for every $v\ge-\log e'$. It contains
$\min(\bar e_0,1/3,\exp(-V_1))>0$ and is downward closed.
Fix
\[
 \bar e_1=\tfrac12\sup\mathcal E_1>0.
\]
There is an element of $\mathcal E_1$ larger than $\bar e_1$;
downward closure therefore gives $\bar e_1\in\mathcal E_1$.
Thus all terminal-scale conditions hold for every
$v\ge-\log \bar e_1$.

\subsection{The physical charge interval}

Fix the reference mass and hence $h\ge1$. Set
\begin{equation}\label{eq:constants-physical-charge}
 e_G=
 \begin{cases}
 \bar e_1,&h=1,\\
 \min\{\bar e_1/\sqrt h,\,1/(\sqrt h\log h)\},&h>1.
 \end{cases}
\end{equation}
This is positive. For $0<e\le e_G$, every running charge
up to level $K$ is at most $\bar e_1$, and
$e\sqrt h\log h\le1$.
The scale proposition therefore applies on the top and
reference tori. The terminal inventory gives
\exteq{43}{TerminalPressure} at their respective stops,
and \eqref{eq:transfer-W3} holds on the reference torus.
The common pair of Theorem~\extref{D}{CoverTuning}
is defined for $N\ge N_{\min}(R_{\mathrm{top}},e)$.
At lower cutoffs only its mass shift is assigned zero.
For $N\ge N_{\mathrm{ref}}(e)$ the same pair is
controlled on the reference torus. Its mass shift is real,
independent of the cutoff volume, and has absolute value
less than $1/8$.
\end{proof}

\subsection{The probe-dependent charge interval}

\begin{proposition}\label{prop:ordered-probe-constants}
Fix the constants of \cref{prop:ordered-constants} and a compact
real probe $\beta$ with $b(dd^*\beta)>0$. Let $C_U$ be the constant
in \exteq{59}{VolumeTransfer}, depending only on the fixed construction,
reference mass and probe, and set $C_f=8C_U+3C_U^2$.
There is a positive $e_\beta\le e_G$ such that, for every
$0<e\le e_\beta$ and $N\ge N_{\rm ref}(e)$,
\eqref{eq:transfer-small-Wick} holds with $\sigma>0$ and
$C_fe^\sigma\le b(dd^*\beta)/2$.
\end{proposition}

\begin{proof} Fix a compact real probe $\beta$. Its buffered support
diameter, energy and uniform rescaled norms are finite by
\exteq{59}{ReferenceResponse}.
The differentiated inverse and interaction constants have
already been fixed; these probe constants are chosen last.
Proposition~\extref{11.1}{ReferenceResponse}
gives $e_\sharp(\beta)>0$ for which all its factor,
one-cell, tree, jet and placement conditions hold for
every $0<e\le e_\sharp(\beta)$, uniformly in
$N\ge N_{\mathrm{ref}}(e)$ and in the allowed placements.
Its proof uses $a_*\in\{h/L,h/L^2\}$ and
$v_*=v_K+O(1)$: polynomial factors have fixed degrees,
the separated-cell errors decay exponentially in $v_*^r$,
and the complement decays faster than the Cauchy-radius
loss. This gives one interval of charges for all cutoffs.

For $0<e\le\min(e_G,e_\sharp(\beta))$,
\exteq{5}{TerminalPressure} and
\exteq{119}{ReferenceResponse} supply the
two inputs in \eqref{eq:transfer-terminal-input}.
The probe placement inequalities imply
\exteq{12}{ReferenceResponse} for $\eta\le1$.
The bounds for the free coefficients and $C_U$ in
\exteq{59}{VolumeTransfer} are finite and independent of
$e,N,T$ and placement. They depend only on the fixed
construction and probe.

Suppose now that $b(dd^*\beta)>0$ and put
$C_f=8C_U+3C_U^2$.
Choose
\begin{equation}\label{eq:constants-probe-charge}
 0<e_\beta\le
 \min\{e_G,e_\sharp(\beta),1,
                         (b(dd^*\beta)/(2C_f))^{1/\sigma}\}.
\end{equation}
The minimum is positive, since $\sigma>0$ and $C_f>0$.
\Cref{prop:transfer-UV-inputs} holds throughout
$(0,e_\beta]$ and gives the finite-cutoff hypotheses of
Theorem~\extref{A}{VolumeTransfer}. Its component clause will be
applied after the component Gaussian bound has been established.
In particular $C_fe^\sigma\le b(dd^*\beta)/2$.
The construction constants were chosen before $h$ and
the probe; the physical charge interval precedes the
probe-dependent restriction.
\end{proof}


\section{Transfer from the reference volumes}\label{sec:uniform-transfer}

The volume-transfer theorem of \cite[Theorem~A]{VolumeTransfer}
controls products of translated Wick polynomials on the cutoff
tori. We verify its hypotheses for the construction chosen in
\cref{sec:ordered-constants}. The partition comparison, the
nonzero-flux estimate and the conditional-moment argument are
proved in that companion.

The selected constants have $r=2$, $s=1/2$, $M=L^k$ with
$k\ge2$, and $0<\varepsilon\le1/200$. The first inequality
of \eqref{eq:constants-L} implies $L\ge13$.
With $h$ as in \eqref{eq:cutoff-scales}, put
\begin{equation}\label{eq:transfer-scale-choice}
 e_K=e\sqrt h,\qquad v_K=-\log e_K,\qquad
 r_K=\min\{L^i:i\in\Z_{\ge0},\ L^i\ge v_K^2\},
 \qquad R^*=hMr_K/L=hL^{T^*}.
\end{equation}
Write $n_*=R^*/\eta$, $N_*=8L^3$ and
$a_L=(L^3-1)/(L-1)^3$. The positive exponent is
\begin{equation}\label{eq:transfer-charge-exponents}
 \mathfrak e=\tfrac12(1/4-9\varepsilon)-12\varepsilon,
 \qquad \sigma=\mathfrak e-2a_L/N_*>0.
\end{equation}
The terminal inventory in \cref{sec:ordered-constants} gives,
for every $0<e\le e_G$, $e\le1/3$, $v_K\ge8$,
$e_K\log h\le1$, and
\begin{equation}\label{eq:transfer-tail-condition}
 6N_*Mv_K^2\exp(-v_K)
             \bigl(v_K\exp(-v_K)+\tfrac12\bigr)\le\pi^2.
\end{equation}
These are the scale hypotheses of the cited theorem. In particular,
\cite[(41), (49)--(51)]{VolumeTransfer} gives
\begin{equation}\label{eq:transfer-W3}
 6N_*e^2R^*\log(1/e)\le\pi^2,
\end{equation}
and
\begin{equation}\label{eq:transfer-volume-error}
 \bar mR^*>Mr_K/L^2\ge r_K\ge v_K,
 \qquad \exp(-\bar mR^*)\le e\sqrt h.
\end{equation}

Fix the compact real probe $\beta$ and $f=dd^*\beta$.
Let $\delta_{\eta,\mu}$ be the centered difference of step
$\eta$ and $J_\eta$ the cube average. Define
\[
 a'_{\eta,\nu}=-\sum_{\mu\ne\nu}\delta_{\eta,\mu}\beta_{\mu\nu},
 \qquad
 f^\eta_{\mu\nu}=\delta_{\eta,\mu}a'_{\eta,\nu}
                       -\delta_{\eta,\nu}a'_{\eta,\mu}.
\]
The covering source is
$\widehat a^o_\nu(x)=\eta^{1/2}(J_\eta a'_{\eta,\nu})(\eta x+\eta e_\nu/2)$.
It satisfies $d\widehat a^o=(f^\eta)_{[\eta]}$, with the
test normalization in \eqref{eq:limit-test-cover}.
Its reflected, periodized copies $\widehat a_c$ and the
variance $v_\eta=\|d\widehat a^o\|^2$ are those of
\cite[Section~2]{ReferenceResponse}.
Let $D_\beta$ be the side of a least coordinate cube containing
$\supp\beta$, and let $Q_\beta$, centered at $z_\beta$, be
that cube enlarged by two in each direction. Its side is
$D_\beta^{\rm buf}=D_\beta+4$. A centered placement is
$v\in\eta\Z^3$ with
$|v+z_\beta-(\ell_c/2)(1,1,1)|_\infty\le\eta/2$,
where $\ell_c=R^*/(2L)$. These are the conventions of
\cite[Section~2]{ReferenceResponse}. Translate $\widehat a^o$
by $v$ before reflecting and periodizing it to form the cochains
$\widehat a_c$. For such a placement,
Proposition~\extref{11.1}{ReferenceResponse} gives, throughout
$0<e\le e_\sharp(\beta)$,
\[
 D_\beta^{\rm buf}+3\eta\le\ell_c,
 \qquad v+Q_\beta\subset[\ell_c/8,7\ell_c/8]^3,
 \qquad \ell_c=R^*/(2L).
\]
Indeed its placement conditions \exteq{121}{ReferenceResponse}
give both inequalities for $\eta\le1$ and a centered placement.
Thus the full support hypothesis \exteq{6}{VolumeTransfer} holds.

For the same placed probe define the ordinary free derivatives by
\[
 N_0(t)=\frac{\det(D^{\rm ch}[q(t/e)\widehat a_0]+\bar m\eta)}
                   {\det(D^{\rm ch}[0]+\bar m\eta)},
 \qquad c_m=N_0^{(m)}(0),\quad m=2,4.
\]
Write
\[
 H_2(x;v)=x^2-v,\qquad H_4(x;v)=x^4-6vx^2+3v^2,
 \qquad W_m(x;v,c)=e^{-m}H_m(x;v)-c.
\]
The reference response uses the interacting mass
$\boldsymbol\mu_N=\bar m\eta+m_{0,N}(e)$, which can have
either sign; this free ratio uses the positive mass $\bar m\eta$.
Let $\E_{n_*}^{(0)}$ be the all-flux law conditioned on zero
flux, at this same interacting mass, and put
\[
 \Theta_m^{(0)}=
 \left|\E_{n_*}^{(0)}\prod_{c\in\{0,\ldots,2L-1\}^3}
 W_m(\langle d\widehat a_c,u\rangle;v_\eta,c_m)\right|^{1/N_*}.
\]
The exact identity \exteq{53}{ReferenceResponse} identifies
this expression with its normalized differentiated response.
The zero-flux law here and the all-flux cutoff law therefore
have exactly the conventions of \cite{VolumeTransfer}.

Let $z_{r,n}$ denote the twist-averaged zero-flux matter factor
\exteq{1}{TerminalPressure}, with mass $\boldsymbol\mu_N$.
For $N\ge N_{\rm ref}(e)$ and
$0<e\le\min\{e_G,e_\sharp(\beta)\}$, the common-pressure
comparison \exteq{5}{TerminalPressure} and the response
bound \exteq{119}{ReferenceResponse} give fixed constants
$C_{\rm mat},C_{\rm ref}<\infty$ such that
\begin{align}
 L^3\log z_{r,n_*}-\log z_{r,Ln_*}&\le C_{\rm mat},\notag\\
 \Theta_4^{(0)}&\le C_{\rm ref}e^{\mathfrak e},\qquad
 (1+|c_2|)\Theta_2^{(0)}\le C_{\rm ref}e^{\mathfrak e}.
                                      \label{eq:transfer-terminal-input}
\end{align}
Both reference tori use the pair selected on $LR^*$;
their common bulk pressure cancels in this difference.
Theorem~\extref{B}{PositiveProbe} and
Proposition~\extref{5.2}{PositiveProbe}, together with
\eqref{eq:transfer-volume-error}, allow $C_{\rm ref}$ to be
enlarged once so that
\begin{equation}\label{eq:transfer-free-input}
 |c_2|+|c_4|^{1/2}\le C_{\rm ref},\qquad
 |c_4-3c_2^2-b(f)|\le C_{\rm ref}(e+\eta^{1/4}).
\end{equation}
Here the support fit gives $R^*\ge D_\beta+1$.
All constants are independent of charge, mesh and placement.
Choose $C_U$ to be the explicit constant
\exteq{59}{VolumeTransfer}, with these reference constants.

\begin{proposition}\label{prop:transfer-UV-inputs}
For the construction and probe above, throughout
$0<e\le\min\{e_G,e_\sharp(\beta)\}$ and
$N\ge N_{\rm ref}(e)$, there are $u_2,u_4\ge0$, independent
of the outer cutoff $T$, such that
\begin{equation}\label{eq:transfer-small-Wick}
 u_4\le C_Ue^\sigma,\qquad
 (1+|c_2|)u_2\le C_Ue^\sigma,\qquad
 |c_4-3c_2^2-b(f)|\le C_U(e^\sigma+\eta^{1/4}).
\end{equation}
For every $T\ge T^*+1$ and finite
$Y\subset(R^*/L)\Z^3$ of points incongruent modulo $R\Z^3$,
\begin{equation}\label{eq:transfer-translated-Wick}
 \left|\E_{N,T}\prod_{y\in Y}
 W_m(F^\sharp(\tau_yf^\eta);v_\eta,c_m)\right|
                      \le u_m^{|Y|},\qquad m=2,4.
\end{equation}
The fixed smooth test satisfies
\begin{equation}\label{eq:transfer-smooth-cumulant}
 |\kappa_4(F^\sharp(f))-e^4b(f)|
 \le(8C_U+3C_U^2)e^{4+\sigma}
                   +C_U(\eta^2+e^4\eta^{1/4}).
\end{equation}
The tests $f^\eta$ have a common compact support,
$f^\eta\to f$ in $L^2$, and $v_\eta\to\|f\|_2^2$.
\end{proposition}

\begin{proof}
The hypotheses of Theorem~\extref{A}{VolumeTransfer}
have just been checked. Its equations (8)--(10), or
Proposition~\extref{6.1}{VolumeTransfer}, give the assertions.
The carrying-torus condition is satisfied because
$n=8L^{T-T^*}n_*$ is a multiple of $p_{\rm cell}=n_*/L$
and is at least $Ln_*$ when $T\ge T^*+1$.
Although $n_*$ is odd, the reference array has $2L$ cells
per side; its half-integer cell widths are precisely those
used in the cited theorem. Its source identity, test-error
bound and variance convergence are \cite[Section~2.2, (19)--(20)]{VolumeTransfer}.
\end{proof}

The translated-product estimate is the input needed for
components: a bound on the parent's fourth cumulant alone
does not bound the fourth cumulants of its components.
Lemma~\extref{8.1}{VolumeTransfer} bounds the invariant
conditional expectations from these products.
Proposition~\extref{8.2}{VolumeTransfer} then gives the
component cumulant estimate. Theorem~\extref{A}{VolumeTransfer}
performs both steps for any prescribed volume-first parent
whose components satisfy the stated Gaussian bound.
We apply this conclusion to the parent and component set
in the proof of \cref{thm:main}.
\section{Tensor Schwinger functions}\label{sec:tensor-axioms}

The Gaussian estimate of \cite[Proposition~6.1]{ContinuumLimits} defines a continuous map
$\Psi\mapsto Q_\Psi$ from $\Sspace_n$ to $L^2(\lambda)$.
To use law-level reflection positivity for arbitrary tensor tests,
we must show that support in a half-space makes this random polynomial
measurable with respect to the corresponding half-space sigma algebra. The Hermite
expansion used for the growth bound does not preserve supports.
We therefore use a compact Fourier expansion for this step.

\subsection{Products and supported tensor tests}

For $\Psi\in\Sspace_k$ and $\Phi\in\Sspace_l$, the tensor evaluation
obeys
\begin{equation}\label{eq:tensor-product-rule}
 Q_\Psi Q_\Phi=Q_{\Psi\otimes\Phi}\quad\text{in }L^1(\lambda).
\end{equation}
Both sides are separately continuous bilinear maps to $L^1$:
for the left side use the $L^2$ bounds and Cauchy--Schwarz, and for
the right side use continuity of the tensor product in Schwartz
seminorms. They agree for elementary tensors. Their linear span
is dense, so approximation in each variable proves
\eqref{eq:tensor-product-rule}. The same argument gives
\begin{equation}\label{eq:tensor-equivariance}
 Q_{\overline\Psi}=\overline{Q_\Psi},\qquad
 Q_{g\cdot\Psi}=\widehat g Q_\Psi
 \quad\text{if }g_*\lambda=\lambda.
\end{equation}
Here the second identity uses the exterior-square action, including
the tensor indices.

We recall the needed supported expansion. Let $O\subset\R^3$ be open
and $\Psi\in C_c^\infty(O^n)$. Its support projects into compact
sets $K_a\subset O$. Choose $\chi_a\in C_c^\infty(O)$ equal to one
on $K_a$, and a cube of side $s$ containing all these supports in
its interior. Extend $\Psi$ periodically from this cube in all its
variables. Its Fourier coefficients $c_{\mathbf k}$ obey, for
every integer $p$,
\[
 |c_{\mathbf k}|\le C_{p,s,\Psi}(1+|\mathbf k|^2)^{-p},
 \qquad \mathbf k=(k_1,\ldots,k_n)\in(\Z^3)^n.
\]
This follows by applying $(1-\Delta)^p$ and integrating by parts;
there are no boundary terms. Multiplication of the Fourier series by
$\prod_a\chi_a(y_a)$ gives the identity on all space
\begin{equation}\label{eq:supported-tensor-expansion}
 \Psi=\sum_{\mathbf k}c_{\mathbf k}
       \bigotimes_{a=1}^n
       \left(\chi_a(y_a)e^{2\pi i k_a\cdot y_a/s}\right).
\end{equation}
The seminorm $q_m$ of each tensor factor product grows at most as
$(1+|\mathbf k|^2)^{m/2}$. Taking $p=m+3n$ makes the series
absolutely convergent in that seminorm. The same argument applies
componentwise to two-form tensor tests. It also proves the density
of elementary tensors used above; arbitrary Schwartz tests are
approximated by growing compact cutoffs.

Each elementary tensor in \eqref{eq:supported-tensor-expansion}
has factors supported in $O$, so its random evaluation is
$\Sigma_O$-measurable and belongs to $L^2(\lambda)$ by the moment
bound. Continuity of $Q$ and closedness of $L^2(\Sigma_O,\lambda)$
give
\begin{equation}\label{eq:tensor-local-measurability}
 Q_\Psi\in L^2(\Sigma_O,\lambda)
       \qquad(\Psi\in C_c^\infty(O^n)).
\end{equation}

This extends to tests supported in the closed half-space
$H=\{x:v\cdot x\ge c\}$, with $|v|=1$. Choose a fixed compact
cutoff $\chi$ equal to one near zero and put
\[
 \Psi_j(y)=\prod_{a=1}^n\chi(y_a/j)\,
           \Psi(y_1-v/j,\ldots,y_n-v/j).
\]
Each $\Psi_j$ has compact support in
$O^n$, $O=\{x:v\cdot x>c\}$. Translation is continuous in Schwartz
space, and the growing cutoff tends to one in its action on each
Schwartz function. More explicitly, Leibniz' rule gives
\[
 q_m(\Psi_j-\Psi)\le C_{m,n,\chi}j^{-1}q_{m+1}(\Psi).
\]
Consequently $Q_{\Psi_j}\to Q_\Psi$ in $L^2(\lambda)$, and
\begin{equation}\label{eq:tensor-closed-halfspace}
 Q_\Psi\in L^2(\Sigma_O,\lambda)\qquad
                   (\supp\Psi\subset H^n).
\end{equation}
Thus tests meeting the boundary plane belong to the same
positive-side $L^2$ space.

\subsection{The five conditions}

\begin{proposition}\label{prop:tensor-OS}
Let $\lambda$ be a Gaussian-dominated law of real tempered two-forms.
Suppose it is invariant under Euclidean motions and reflection
positive at every plane through the origin. Then its tensor
Schwinger functions satisfy \textup{(E0$'$)}--\textup{(E3)} of
\cref{thm:main}. If, for every $A,B\in L^2(\lambda)$,
\[
 \E_\lambda[\overline A\,U(a)B]\longrightarrow
       \overline{\E_\lambda A}\,\E_\lambda B
       \quad\text{as }|a|\longrightarrow\infty,
\]
they also satisfy \textup{(E4)}.
\end{proposition}

\begin{proof}
The Hermite estimate of \cite[Proposition~6.1]{ContinuumLimits} proves \textup{(E0$'$)}. Covariance follows from
\eqref{eq:tensor-equivariance} and invariance of expectation.
Permutation symmetry holds on elementary tensors because the random
field evaluations commute; continuity gives \textup{(E3)} on all
tensor tests.

For \textup{(E2)}, let $r$ be the stated reflection and $H$ the
chosen closed half-space. By \eqref{eq:tensor-closed-halfspace},
$A=\sum_{k=0}^MQ_{\Psi_k}$ lies in its positive-side $L^2$ space,
where $Q_{\Psi_0}=\Psi_0$. Reflection positivity extends from
bounded characters to this whole space by their density and the
boundedness of the reflection form. Equations
\eqref{eq:tensor-product-rule}--\eqref{eq:tensor-equivariance} give
\[
 \overline{\mathcal R_rQ_{\Psi_k}}
       =Q_{\Theta_k\Psi_k},\qquad
 Q_{\Theta_k\Psi_k}Q_{\Psi_l}
       =Q_{\Theta_k\Psi_k\otimes\Psi_l}\quad\text{in }L^1.
\]
Therefore
\[
 0\le\E_\lambda[(\overline{\mathcal R_rA})A]
   =\sum_{k,l=0}^MS_{k+l}(\Theta_k\Psi_k\otimes\Psi_l).
\]
This is the tensor reflection form with the conjugation and signs
in \cref{thm:main}.

Finally, \eqref{eq:tensor-product-rule} and translation covariance
identify the expression in \textup{(E4)} with
$\E_\lambda[Q_\Psi U(a)Q_\Phi]$. Both random polynomials are in
$L^2(\lambda)$. Apply the stated clustering with
$A=\overline{Q_\Psi}$ and $B=Q_\Phi$ to obtain
$\E_\lambda Q_\Psi\,\E_\lambda Q_\Phi$.
\end{proof}

For the components in \cref{prop:limit-nine-mirrors}, the clustering
hypothesis holds. The covariance theorem in \cref{sec:rotations}
will supply their Euclidean invariance and reflection positivity
at every plane.


\section{Completion of the field-strength construction}
\label{sec:main-theorem-closure}

\begin{proof}[Proof of \cref{thm:main}]
Choose the construction constants by \cref{prop:ordered-constants},
in the stated order.
In particular $\sigma>0$ in
\eqref{eq:transfer-charge-exponents}.
Take the positive physical ceiling $e_G$ in
\eqref{eq:constants-physical-charge}. Theorem~\extref{D}{CoverTuning} supplies real bare mass shifts
$m_{0,N}(e)$, independent of the volume and twist,
for $0<e\le e_G$; below its defining mesh threshold they
are set to zero. Their bound
\exteq{65}{CoverTuning} gives
$|m_{0,N}(e)|<1/8$ for every $N$.
The all-flux laws at \eqref{eq:cutoff-scales} are therefore
defined at every cutoff, even when the total lattice mass
has either sign.
The physical ceiling in \cref{thm:main} is
$e_{\mathrm{phys}}=e_G$.
Theorem~\extref{A}{PositiveProbe} proves the existence
of the probes in \eqref{eq:positive-coefficient}.
Fix any such probe and $0<e\le e_\beta$, with
$e_\beta$ supplied by \cref{prop:ordered-probe-constants}.

For this probe,
Proposition~\extref{11.1}{ReferenceResponse}
provides $e_\sharp(\beta)>0$.
The terminal partition bound
\exteq{5}{TerminalPressure} and the response estimate
\exteq{119}{ReferenceResponse} supply
both parts of \eqref{eq:transfer-terminal-input}.
The free-coefficient bound
\eqref{eq:transfer-free-input} supplies its remaining
reference input. Thus \cref{prop:transfer-UV-inputs}
applies throughout this interval, with one $C_U$
independent of the cutoff, volume, placement and limit law.
Its estimates hold whenever
$N\ge N_{\rm ref}(e)$ and $T\ge T^*(e)+1$.
These are finite thresholds for each fixed $e$.

\emph{Volume-first laws and their Euclidean components.}
The tightness estimate \eqref{eq:limit-uniform-tightness}
shows that every fixed-spacing volume-limit set
$\mathcal V_N$ is nonempty and that the collection of
these limits is tight on $H_{-4}$.
Consequently volume-first limit points exist.
Let $\mu$ be any one of them, and write
$\mu^{(N_i)}\to\mu$, where $N_i\to\infty$ and
$\mu^{(N_i)}\in\mathcal V_{N_i}$.
For each $i$, choose a volume index $T_i$ along a
defining subsequence for $\mu^{(N_i)}$ such that
\[
 T_i\ge i,\qquad
 d_{\rm w}(P_{N_i,T_i},\mu^{(N_i)})\le1/i,
\]
where $d_{\rm w}$ metrizes weak convergence on $H_{-4}$.
Then
$d_{\rm w}(P_{N_i,T_i},\mu)
\le1/i+d_{\rm w}(\mu^{(N_i)},\mu)\to0$.
This is the joint sequence in
\cref{prop:limit-nine-mirrors};
both indices diverge, and all fixed finite-family
moments converge along it.

\Cref{prop:limit-nine-mirrors} gives closedness,
translation invariance, Gaussian domination with constant
one, and invariance and reflection positivity at all
planes in the nine mirror families.
Sections~\extref{5--6}{ErgodicComponents} transfer
these properties to a measurable set of full measure
in the translation-ergodic decomposition.
The Gaussian bound holds there for all real Schwartz
tests at once, by the countable-density and Fatou
argument in \exteq{20}{ErgodicComponents}.
Section~\extref{5}{ErgodicComponents} transfers each additional mirror
by transporting the conditional laws and using
uniqueness of the decomposition. Intersecting the
nine full-measure sets gives all mirror forms in
each component; these reflections also generate
the cubic group.

\Cref{thm:rotation-covariance} now applies
separately to $\mu$ and to every component in this
full-measure set. It proves their invariance under
the full Euclidean group and reflection positivity
at every plane.
For an ergodic component, the sharp cone and
\exteq{15}{ErgodicComponents} identify the
translation-invariant projection with the constant
projection. Proposition~\extref{4.1}{ErgodicComponents}
therefore gives clustering for all $L^2$ observables
as $|a|\to\infty$.
\Cref{prop:tensor-OS} gives
\textup{(E0$'$)}--\textup{(E3)} for the parent law and
\textup{(E0$'$)}--\textup{(E4)} for these components,
with exactly the tensor action and growth bound of
\cref{thm:main}.
In particular, supported tensor evaluations belong
to the positive-side $L^2$ space even when their
supports meet the reflection plane, by
\eqref{eq:tensor-closed-halfspace}.
For clustering the relevant identity is
\[
 S_{k+l}^\xi(\Psi\otimes\tau_a^{\otimes l}\Phi)
   =\E_{\mu_\xi}[Q_\Psi\,U(a)Q_\Phi],
                 \qquad Q_\Psi,Q_\Phi\in L^2(\mu_\xi).
\]
It gives \eqref{eq:main-cluster} on all tensor Schwartz
tests, using \(A=\overline{Q_\Psi}\) in the
$L^2$ clustering formula.

\emph{Fourth cumulants in the same components.}
The prescribed parent $\mu$ is the volume-first limit of the
same all-flux laws, at the same volume-independent mass assignment,
used in \cref{prop:transfer-UV-inputs}. Their physical interpolation
and weak topology are those of Theorem~\extref{A}{VolumeTransfer}.
The preceding argument gives translation invariance of $\mu$
and Gaussian domination with constant one on a full set of
its translation-ergodic components. Thus the component hypothesis
\exteq{11}{VolumeTransfer} holds for this parent.

The reference, placement and charge hypotheses of that theorem
were verified in \cref{sec:uniform-transfer}. Its conclusion
\exteq{12}{VolumeTransfer} gives a full-measure component set on which
\[
 |\kappa_4^{\mu_\xi}(f)-e^4b(f)|\le C_fe^{4+\sigma},
 \qquad C_f=8C_U+3C_U^2.
\]
The prescribed-parent clause keeps the law $\mu$ fixed;
the coefficient subsequence used in its proof does not select
a different parent. Intersect this set with the full-measure
set already used for the tensor OS conditions.
The charge choice \eqref{eq:constants-probe-charge} gives
$C_fe^\sigma\le b(f)/2$, hence
$\kappa_4^{\mu_\xi}(f)\ge e^4b(f)/2>0$ on this intersection.
The constant is independent of the parent and component.
This proves \eqref{eq:main-cumulant}; the full-measure set
is asserted for the fixed charge, probe and parent law.
A centered Gaussian variable has zero fourth cumulant, so
these component laws are non-Gaussian.
\end{proof}
\section{Rotational covariance}\label{sec:rotations}

The rotation argument starts with a field insertion separated from
the two observables in a reflection form. Gaussian domination bounds
the insertion in the reconstructed Hilbert space, independently of
the number of factors in either observable. The spectral cone then
allows complex translations. Together these bounds give bounded
operators for insertions concentrated at a point of the
$(x_1,x_2)$-plane. Closedness removes the extra angular frequencies
at two points. At higher orders, reflection positivity gives a positive
kernel for the highest angular coefficient. Polygon divisibility and a
convexity estimate then reduce every order to the two-point result.


We use the Gaussian moment bounds, reconstruction and spectral cone
recalled in \cref{sec:probability}. The parent law and its components
have these properties by \cref{prop:limit-nine-mirrors}. Their
closedness enters the two-point frequency calculation and the
removal of contact terms.
Let $\mu$ be a translation-invariant, Gaussian-dominated law,
invariant and reflection positive at $x_1=0$.
We first work with this reconstruction, writing
$\mathcal H=\mathcal H_1$, $H=H_1$, and
$U_t=U(te_1)$.
Let $\mathcal A_+$ denote the positive-side measurable functions
belonging to every $L^p(\mu)$, $p<\infty$. It is an algebra.
Their classes are dense in $\mathcal H$, since $\mathcal A_+$
contains the bounded positive-side functions.

Use the Hodge isometry
$*f=(f_{23},-f_{13},f_{12})$ to write
$\mathbb B(b)=F(*^{-1}b)$, complex linearly. The completion recalled
in \cref{sec:probability} extends this map to vector-valued $L^2$ tests.
The action on these tests is
\[
 (g\cdot b)(x)=\det(g)\,g\,b(g^{-1}x)
 \quad(g\in O(3)),\qquad
 (\tau_a b)(x)=b(x-a).
\]
In particular
$\Theta_1\mathbb B(b)=\mathbb B(\overline{r_1\cdot b})$
and $\Theta_1U_t=U_{-t}\Theta_1$.
Write $L^2_{[a,b]}$ for the vector tests vanishing almost everywhere
outside $a\le x_1\le b$.

\subsection{The insertion estimate}\label{subsec:rotation-insertion}

For $\ell>0$ and $h\in L^2_{[0,\ell]}$, consider
$M_\ell(h)\Phi=e^{\mathbb B(h)}U_\ell\Phi$ for bounded
positive-side $\Phi$. This function lies in $\mathcal A_+$.
We first prove
\begin{equation}\label{eq:rotation-source-bound}
 \|[M_\ell(h)\Phi]_1\|
       \le e^{\|\Re h\|_2^2/2}\|[\Phi]_1\|.
\end{equation}

Set $h^\#=\tau_{\ell e_1}(\overline{r_1\cdot h}+h)$.
Reflection and translation invariance give
\begin{equation}\label{eq:rotation-source-doubling}
 \|[M_\ell(h)\Phi]_1\|^2
   =\langle[\Phi]_1,[M_{2\ell}(h^\#)\Phi]_1\rangle.
\end{equation}
The two summands in $h^\#$ have disjoint slab interiors, and
$\tau_{2\ell e_1}\overline{r_1\cdot h^\#}=h^\#$.

For a source $h'$ on $[0,\ell']$ with
$\tau_{\ell'e_1}\overline{r_1\cdot h'}=h'$, let
$h'^{(n)}=\sum_{j=0}^{n-1}\tau_{j\ell'e_1}h'$.
Doubling $h'^{(n)}$ gives $h'^{(2n)}$. Thus, with
$m_a=\|[M_{2^a\ell'}(h'^{(2^a)})\Phi]_1\|$,
\[
 m_a^2\le\|[\Phi]_1\|m_{a+1},\qquad
 m_0\le\|[\Phi]_1\|^{1-2^{-a}}m_a^{2^{-a}}.
\]
On the other hand, the integrand in the doubling identity is bounded
in modulus by
\[
 \|\Phi\|_\infty^2e^{\mathbb B(\Re h'^{(2^{a+1})})}.
\]
The real translates have disjoint interiors, so Gaussian domination
gives
\[
 m_a^2\le\|\Phi\|_\infty^2
       e^{\,2^a\|\Re h'\|_2^2}.
\]
Taking the indicated powers and sending $a\to\infty$ proves
\eqref{eq:rotation-source-bound} for symmetric sources, including
$\|[\Phi]_1\|=0$. Apply this to $h^\#$, whose real part has squared
norm $2\|\Re h\|_2^2$, in
\eqref{eq:rotation-source-doubling}. Cauchy--Schwarz proves the
bound for arbitrary $h$.

\begin{proposition}\label{prop:rotation-insertion}
For $\ell>0$ and $f\in L^2_{[0,\ell]}$ there is a unique bounded
operator $\mathcal B_\ell(f)$ on $\mathcal H$ satisfying
\[
 \mathcal B_\ell(f)[\Phi]_1
       =[\mathbb B(f)U_\ell\Phi]_1\quad(\Phi\in\mathcal A_+).
\]
It is complex linear in $f$, and
\begin{equation}\label{eq:rotation-insertion-bound}
 \|\mathcal B_\ell(f)\|\le k_1\|f\|_2,\qquad
 k_1=\sqrt{2\exp(1)}.
\end{equation}
For real $f$, the constant is $\exp(1/2)$.
If $f\in L^2_{[c,c']}$ with $0\le c<c'\le\ell$, then
\begin{equation}\label{eq:rotation-insertion-local}
 \mathcal B_\ell(f)=e^{-cH}
       \mathcal B_{c'-c}(\tau_{-ce_1}f)e^{-(\ell-c')H}.
\end{equation}
For $f\in L^2_{[-d,d]}$, set
$\mathcal B_d^{\mathrm c}(f)=\mathcal B_{2d}(\tau_{de_1}f)$.
Its matrix elements and adjoint are
\begin{align}
 \langle[\Phi']_1,\mathcal B_d^{\mathrm c}(f)[\Phi]_1\rangle
   &=\E_\mu[(\Theta_1U_d\Phi')\mathbb B(f)U_d\Phi],
                         \label{eq:rotation-insertion-centered}\\
 \mathcal B_d^{\mathrm c}(f)^*
   &=\mathcal B_d^{\mathrm c}(\overline{r_1\cdot f}).\notag
\end{align}
For $r\ge|t|$ and $a\in\R$, one also has
\begin{multline}\label{eq:rotation-insertion-shift}
 e^{-(r+t)H+ia\mathsf P_2}\mathcal B_d^{\mathrm c}(f)
           e^{-(r-t)H-ia\mathsf P_2}\\
       =\mathcal B_{d+r}^{\mathrm c}(\tau_{te_1+ae_2}f).
\end{multline}
\end{proposition}

\begin{proof}
For real $f$ and bounded $\Phi$, the map
$z\mapsto[e^{z\mathbb B(f)}U_\ell\Phi]_1$ is entire.
Indeed its exponential series converges locally uniformly in
$L^2(\mu)$, dominated on $|z|\le R$ by
$\|\Phi\|_\infty e^{R|\mathbb B(f)|}$.
The quotient map is a contraction. On $|z|=\|f\|_2^{-1}$,
\eqref{eq:rotation-source-bound} bounds its norm by
$\exp(1/2)\|[\Phi]_1\|$. Cauchy's coefficient estimate gives the
real insertion bound. Applying it to $\Re f$ and $\Im f$, and using
$\|\Re f\|_2+\|\Im f\|_2\le\sqrt2\|f\|_2$, gives
\eqref{eq:rotation-insertion-bound}.
The estimate sends null classes to null classes, so it defines the
operator on the dense classes of bounded functions and extends it
uniquely to $\mathcal H$.

For $\Phi\in\mathcal A_+$, its bounded truncations converge in $L^4$.
Since $\mathbb B(f)\in L^4$, their products with the inserted field
converge in $L^2$ by H\"older's inequality. This proves the asserted
formula on $\mathcal A_+$ and complex linearity in $f$.

Equation \eqref{eq:rotation-insertion-local} follows by applying
$U_c$ to
$\mathbb B(\tau_{-ce_1}f)U_{\ell-c}\Phi$ and using
\exteq{6}{ErgodicComponents}. Translation by $-de_1$ in the reflection
form gives \eqref{eq:rotation-insertion-centered}. Applying reflection
and conjugation to that integrand gives the adjoint formula.
For \eqref{eq:rotation-insertion-shift}, apply the two exterior
semigroups to a class $[\Phi]_1$. The resulting function is
\[
 \mathbb B(\tau_{(d+r)e_1+te_1+ae_2}f)\,
       U_{2(d+r)}\Phi.
\]
The test $\tau_{te_1+ae_2}f$ is supported in $[-d-r,d+r]$,
since $|t|\le r$. The displayed function therefore defines its centered
insertion at $d+r$, which is the right side. Equality on these dense
classes proves the operator identity.
\end{proof}

In particular, for $\Phi_1,\Phi_2\in\mathcal A_+$ measurable on
$x_1>d$,
\begin{equation}\label{eq:rotation-separated-insertion}
 |\E_\mu[(\Theta_1\Phi_1)\mathbb B(f)\Phi_2]|
 \le k_1\|f\|_2\,\|[U_{-d}\Phi_1]_1\|
                         \|[U_{-d}\Phi_2]_1\|.
\end{equation}
This is \eqref{eq:rotation-insertion-centered} applied to the two
translated observables. The constant does not depend on their degrees.

\subsection{Complex translations and point plates}
\label{subsec:rotation-point-plates}

From now on assume also the sharp cone
$\mathsf E_1\{(E,p_\perp):E<|p_2|\}=0$.
For the laws under construction, Proposition~\extref{3.1}{ErgodicComponents} supplies this from the diagonal mirror.

For $\xi=(\xi^0,\xi^1)\in\C^2$, put
\[
 d_{\mathcal W}(\xi)=\Re\xi^0-|\Im\xi^1|,\qquad
 \mathcal W=\{d_{\mathcal W}>0\},\qquad
 \xi^*=(\overline{\xi^0},-\overline{\xi^1}),\quad n_0=(1,0).
\]
The cone \exteq{8}{ErgodicComponents} gives contractions
\[
 T(\xi)=\int e^{-\xi^0E+i\xi^1p_2}\,\dd\mathsf E_1(E,p_\perp)
                  \quad(\xi\in\overline{\mathcal W}).
\]
They obey $T(\xi)T(\xi')=T(\xi+\xi')$,
$T(\xi)^*=T(\xi^*)$, and $T(\xi)\Omega_1=\Omega_1$.
They are strongly continuous on the closed tube by spectral
dominated convergence. In its interior they are holomorphic in
operator norm: on $d_{\mathcal W}\ge\epsilon>0$ the first and
second derivatives of the multiplier are bounded by constants
times $E e^{-\epsilon E}$ and $E^2e^{-\epsilon E}$.
Taylor's formula therefore gives a uniform quadratic remainder.
For real $y=(y^0,y^1)$ with $y^0\ge0$,
$T(y)[\Phi]_1=[U(y^0e_1+y^1e_2)\Phi]_1$.

Write $x=(y,x_3)$ with $y=(x_1,x_2)$. Set
\begin{gather*}
 e^+=e_1+ie_2,\quad e^-=e_1-ie_2,\quad e^3=e_3,\\
 w(+)=1,\quad w(-)=-1,\quad w(3)=0,\\
 s(+)=s(-)=1,\quad s(3)=-1.
\end{gather*}
A label $\lambda=(\sigma,\chi)$ has
$\sigma\in\{+,-,3\}$ and real $\chi\in L^2(\R)$.
For $\alpha\in L^2(\R^2;\C)$ its plate is
$(\alpha\otimes\lambda)(x)=\alpha(y)\chi(x_3)e^\sigma$.
Thus $\|\alpha\otimes\lambda\|_2\le\sqrt2\|\alpha\|_2\|\chi\|_2$.
For rotations $R_t$ in the $(x_1,x_2)$-plane and the coordinate
reflection $r_1$, direct multiplication gives
\begin{align}
 R_t\cdot(\alpha\otimes\lambda)
   &=e^{-iw(\sigma)t}(\alpha\circ R_{-t})\otimes\lambda,
                              \label{eq:rotation-plate-actions}\\
 \overline{r_1\cdot(\alpha\otimes\lambda)}
   &=s(\sigma)(\overline\alpha\circ r_1)\otimes\lambda.\notag
\end{align}

\begin{proposition}\label{prop:rotation-point-plates}
For each label $\lambda$ there are bounded operators
$\mathcal B_\lambda(\xi,\xi')$, $\xi,\xi'\in\mathcal W$,
holomorphic in operator norm and real linear in $\chi$, such that
\begin{equation}\label{eq:rotation-point-bound}
 \|\mathcal B_\lambda(\xi,\xi')\|
 \le\frac{k_2\|\chi\|_2}
          {\min(d_{\mathcal W}(\xi),d_{\mathcal W}(\xi'))},
 \qquad k_2=\frac8\pi\exp(1/2).
\end{equation}
For $\kappa,\kappa'\in\overline{\mathcal W}$,
\begin{align}
 T(\kappa)\mathcal B_\lambda(\xi,\xi')T(\kappa')
   &=\mathcal B_\lambda(\kappa+\xi,\xi'+\kappa'),
                                   \label{eq:rotation-point-shift}\\
 \mathcal B_\lambda(\xi,\xi')^*
   &=s(\sigma)\mathcal B_\lambda(\xi'^*,\xi^*).\notag
\end{align}
If $\alpha\in L^2(\R^2)$ is compactly supported in $|y^0|<d$, then
\begin{equation}\label{eq:rotation-point-integral}
 \mathcal B_d^{\mathrm c}(\alpha\otimes\lambda)
   =\int\alpha(y)\mathcal B_\lambda(dn_0+y,dn_0-y)\,\dd y.
\end{equation}
The integral is a Bochner integral in operator norm.
When $\supp\alpha\subset\{|y^0|\le d'\}$, $d'<d$, its norm is at
most $k_2\|\chi\|_2\|\alpha\|_1/(d-d')$.
\end{proposition}

\begin{proof}
Let $j_\delta=(\pi\delta^2)^{-1}\mathbf1_{\{|y|<\delta\}}$.
Set $A_\delta=\mathcal B_\delta^{\mathrm c}(j_\delta\otimes\lambda)$
and, when both tube distances exceed $\delta$, set
\[
 \mathcal B_\lambda^\delta(\xi,\xi')
   =T(\xi-\delta n_0)A_\delta T(\xi'-\delta n_0).
\]
These operators obey the shift and adjoint identities in
\eqref{eq:rotation-point-shift}. Here $j_\delta$ is real and
reflection invariant, so $A_\delta^*=s(\sigma)A_\delta$.
The real translation identity \eqref{eq:rotation-insertion-shift}
also gives, for $r\ge|y^0|$,
\begin{equation}\label{eq:rotation-point-regularized-translation}
 \mathcal B_\lambda^\delta((r+\delta)n_0+y,(r+\delta)n_0-y)
    =\mathcal B_{\delta+r}^{\mathrm c}((\tau_yj_\delta)\otimes\lambda).
\end{equation}

We prove a bound uniform as $\delta\downarrow0$.
Fix $\xi,\xi'$ and set
$d_0=\min(d_{\mathcal W}(\xi),d_{\mathcal W}(\xi'))$,
$r=(d_0-\delta)/2>0$. For unit vectors $\psi,\psi'$,
\[
 h(z)=\langle\psi,\mathcal B_\lambda^\delta(\xi+z,\xi'-z)\psi'\rangle
\]
is holomorphic on $|\Re z^0|<r$, $|\Im z^1|<r$.
Write $z=y+i(t,t')$. Factoring the exterior translations gives
\[
 h(y+i(t,t'))=
 \langle\psi_1,\mathcal B_{\delta+r}^{\mathrm c}
             ((\tau_yj_\delta)\otimes\lambda)\psi'_1\rangle,
\]
where
$\psi_1=T(\xi-(\delta+r)n_0+i(t,t'))^*\psi$ and
$\psi'_1=T(\xi'-(\delta+r)n_0-i(t,t'))\psi'$.
For $|t'|<r$ these are contractions of the original vectors;
they are independent of $y$.
For $\eta\in L^2((-r,r)^2)$, extended by zero, the integral of
$\eta(y)\tau_yj_\delta$ is $\eta*j_\delta$ in $L^2$.
Its support lies in $|y^0|\le r+\delta$.
The insertion bound and Young's inequality therefore give
\[
 \left|\int\eta(y)h(y+i(t,t'))\,\dd y\right|
   \le2\exp(1/2)\|\eta\|_2\|\chi\|_2.
\]
Hence
$\int_{(-r,r)^2}|h(y+i(t,t'))|^2\,\dd y
\le4\exp(1)\|\chi\|_2^2$.
Apply the mean-value inequality on a polydisc of radius $r'<r$
and bound its integral by that over the four-dimensional box.
This yields
\[
 |h(0)|^2\le(\pi r'^2)^{-2}(2r)^2\,4\exp(1)\|\chi\|_2^2.
\]
Letting $r'\uparrow r$ and taking arbitrary unit vectors proves
\begin{equation}\label{eq:rotation-regularized-point-bound}
 \|\mathcal B_\lambda^\delta(\xi,\xi')\|
       \le\frac{k_2\|\chi\|_2}{d_0-\delta}.
\end{equation}

To identify the limit, put $r_0=d_0/4$.
For $\delta<r_0$ the functions
$h_\delta(z)=\langle\psi,\mathcal B_\lambda^\delta
(\xi+z,\xi'-z)\psi'\rangle$ form a bounded normal family on
$|z^0|,|z^1|<r_0$ by
\eqref{eq:rotation-regularized-point-bound}.
For $\eta\in L^2((-r_0,r_0)^2)$, the real translation identity
and an additional equal time translation on both sides give
\begin{multline}\label{eq:rotation-point-limit-test}
 \int\eta(y)h_\delta(y)\,\dd y\\
 =\langle\psi,T(\xi-2r_0n_0)
       \mathcal B_{2r_0}^{\mathrm c}((\eta*j_\delta)\otimes\lambda)
       T(\xi'-2r_0n_0)\psi'\rangle.
\end{multline}
Since $\eta*j_\delta\to\eta$ in $L^2$, this has a unique limit.
Any two holomorphic subsequential limits thus have the same
integrals on the real square; continuity identifies them there,
and the identity theorem identifies them on the polydisc.
Consequently every matrix element at $z=0$ converges.
The bound \eqref{eq:rotation-regularized-point-bound} defines a
bounded limiting sesquilinear form, represented by
$\mathcal B_\lambda(\xi,\xi')$ with
\eqref{eq:rotation-point-bound}. This is weak operator convergence.
The shift, adjoint, and real linearity identities pass to these
limits.

For holomorphy, fix $(\xi_0,\xi'_0)$ and a positive $\epsilon$
smaller than both tube distances. Near this pair the shift identity
factors the operator as
\[
 T(\xi-\xi_0+\epsilon n_0)\,
 \mathcal B_\lambda(\xi_0-\epsilon n_0,\xi'_0-\epsilon n_0)\,
 T(\xi'-\xi'_0+\epsilon n_0).
\]
The two exterior factors are holomorphic in operator norm.

Finally choose $d'<d$ containing the time support of $\alpha$.
Integrating
\eqref{eq:rotation-point-regularized-translation} with $r=d-\delta$
gives
\[
 \int\alpha(y)\mathcal B_\lambda^\delta(dn_0+y,dn_0-y)\,\dd y
       =\mathcal B_d^{\mathrm c}((\alpha*j_\delta)\otimes\lambda).
\]
The right side converges in norm to the left side of
\eqref{eq:rotation-point-integral}. For matrix elements on the
left, \eqref{eq:rotation-regularized-point-bound} gives the
integrable majorant
$2k_2\|\chi\|_2|\alpha(y)|/(d-d')$ for sufficiently small $\delta$.
Dominated convergence identifies the limiting matrix elements.
Holomorphy makes the operator factor norm continuous on the
support. Multiplication by the measurable scalar $\alpha$ gives
a strongly measurable integrand, whose norm is integrable by
\eqref{eq:rotation-point-bound} and $\alpha\in L^1$. This proves
Bochner integrability, the formula and its stated norm bound.
\end{proof}

\subsection{Ordered field kernels}\label{subsec:rotation-ordered-kernels}

The intermediate arguments of a product of point plates matter
only through their sum. Indeed, for $\kappa_1,\kappa_2\in\mathcal W$
and sufficiently small $\epsilon>0$, the shift identity gives
\[
 \mathcal B_\lambda(\xi,\kappa_1)
       \mathcal B_{\lambda'}(\kappa_2,\xi')
 =\mathcal B_\lambda(\xi,\epsilon n_0)
       T(\kappa_1+\kappa_2-2\epsilon n_0)
       \mathcal B_{\lambda'}(\epsilon n_0,\xi').
\]
The same expression results from any other decomposition of
$\kappa_1+\kappa_2$ into two points of $\mathcal W$.
An exterior argument adjacent to the vacuum can likewise be
changed arbitrarily in $\mathcal W$, since $T(\kappa)\Omega_1=
T(\kappa)^*\Omega_1=\Omega_1$.

For $N\ge2$ and $\xi_1,\ldots,\xi_{N-1}\in\mathcal W$, define
\begin{multline}\label{eq:rotation-ordered-kernel}
 G_N(\xi;\lambda)=
 \langle\Omega_1,\mathcal B_{\lambda_1}(\kappa,\xi_1/2)
       \left(\prod_{a=2}^{N-1}
          \mathcal B_{\lambda_a}(\xi_{a-1}/2,\xi_a/2)\right)\\
       \mathcal B_{\lambda_N}(\xi_{N-1}/2,\kappa)\Omega_1\rangle,
       \qquad\kappa\in\mathcal W.
\end{multline}
Set $G_1(\lambda)=\langle\Omega_1,
\mathcal B_\lambda(\kappa,\kappa)\Omega_1\rangle$.
Both definitions are independent of $\kappa$.
For $x_1,x_{a+1}-x_a\in\mathcal W$, define the vector
\begin{multline}\label{eq:rotation-ordered-vector}
 \Psi(x;\lambda)=
 \mathcal B_{\lambda_1}(x_1,(x_2-x_1)/2)
       \left(\prod_{a=2}^{n-1}
          \mathcal B_{\lambda_a}((x_a-x_{a-1})/2,(x_{a+1}-x_a)/2)\right)\\
       \mathcal B_{\lambda_n}((x_n-x_{n-1})/2,\kappa)\Omega_1.
\end{multline}
Products here and below are ordered with increasing index from
left to right; an empty product is the identity.
For $n=1$ this means
$\mathcal B_{\lambda_1}(x_1,\kappa)\Omega_1$; for $n=0$ it means
$\Omega_1$. It too is independent of $\kappa$.
Put $s(\lambda)=\prod_a s(\sigma_a)$, $w(\lambda)=\sum_a w(\sigma_a)$,
and let $\bar\lambda$ exchange the $+$ and $-$ labels.

\begin{proposition}\label{prop:rotation-ordered-kernels}
The kernels and vectors are holomorphic on their stated domains.
For $N\ge2$,
\begin{equation}\label{eq:rotation-ordered-kernel-bound}
 |G_N(\xi;\lambda)|
 \le\prod_{a=1}^N(2k_2\|\chi_a\|_2)
       \left(\min_a d_{\mathcal W}(\xi_a)\right)^{-N}.
\end{equation}
Let $\alpha_1,\ldots,\alpha_N\in L^2(\R^2)$ have compact
time-ordered supports:
$\sup_{\supp\alpha_a}y^0<\inf_{\supp\alpha_{a+1}}y^0$.
Then
\begin{equation}\label{eq:rotation-ordered-density}
 \E_\mu\prod_a\mathbb B(\alpha_a\otimes\lambda_a)
 =\int\prod_a\alpha_a(y_a)\,
       G_N(y_2-y_1,\ldots,y_N-y_{N-1};\lambda)\,\dd Y.
\end{equation}
If the supports are also in $y^0>0$, their product class is
the Bochner integral of $\prod_a\alpha_a(y_a)\Psi(Y;\lambda)$.
One has $G_1=0$,
$\overline{G_N(\xi;\lambda)}=G_N(\bar\xi;\bar\lambda)$, and
\begin{multline}\label{eq:rotation-ordered-Gram}
 \langle\Psi(x;\lambda),\Psi(z;\lambda')\rangle
 =s(\lambda)G_{n+n'}((x_n-x_{n-1})^*,\ldots,(x_2-x_1)^*,\\
       x_1^*+z_1,z_2-z_1,\ldots,z_{n'}-z_{n'-1};
       \lambda_n,\ldots,\lambda_1,\lambda'_1,\ldots,\lambda'_{n'})
\end{multline}
when $n,n'\ge1$. With a vacuum in the first slot the inner
product is $G_{n'}$ on the successive differences of $z$;
with a vacuum in the second slot it is $s(\lambda)G_n$
on the reversed starred differences and reversed labels.
Furthermore
$T(\kappa')\Psi(x;\lambda)=\Psi(x+\kappa';\lambda)$ for
$\kappa'\in\overline{\mathcal W}$, where every point is translated.
\end{proposition}

\begin{proof}
Holomorphy follows from that of the factors. In the definition of
$G_N$, choose $\kappa=tn_0$ with
$t\ge\frac12\min_a d_{\mathcal W}(\xi_a)$ and apply
\eqref{eq:rotation-point-bound} to every factor. This proves
\eqref{eq:rotation-ordered-kernel-bound}.

To prove the density formula, first put the supports in $y^0>0$.
Choose disjoint time slabs $[a_j,b_j]$, in order, with each support
strictly inside its slab, $a_1\ge0$, and $b_j\le a_{j+1}$.
Let $f_j=\alpha_j\otimes\lambda_j$,
$c_j=(a_j+b_j)/2$, $d_j=(b_j-a_j)/2$, and
$I_j=\mathcal B_{d_j}^{\mathrm c}(\tau_{-c_je_1}f_j)$.
The insertion identity and its localization give
\begin{equation}\label{eq:rotation-ordered-product-class}
 \left[\prod_j\mathbb B(f_j)\right]_1
 =T(a_1n_0)I_1T((a_2-b_1)n_0)I_2\cdots I_N\Omega_1.
\end{equation}
For precision, the backward induction uses
$\Phi'=\prod_{j>j_0}\mathbb B(\tau_{-b_{j_0}e_1}f_j)$,
which belongs to $\mathcal A_+$.
The tail product is $U_{b_{j_0}}\Phi'$, so insertion of $f_{j_0}$
is $\mathcal B_{b_{j_0}}(f_{j_0})[\Phi']_1$.
Localization supplies $T(a_{j_0}n_0)I_{j_0}$; the induction
on the translated tail supplies the next nonnegative gap
$a_{j_0+1}-b_{j_0}$ and the remaining factors.

Use \eqref{eq:rotation-point-integral} to write
\[
 I_j=\int\alpha_j(y_j)
       \mathcal B_{\lambda_j}(y_j-a_jn_0,b_jn_0-y_j)\,\dd y_j.
\]
On the compact supports all these operators are norm bounded;
the $\alpha_j$ are integrable. The product in
\eqref{eq:rotation-ordered-product-class} is therefore one Bochner
integral. Absorb each gap semigroup into a neighboring point plate.
The inner arguments between plates $j$ and $j+1$ then add to
$y_{j+1}-y_j$. Replace them by equal halves using the product
identity above. The initial argument becomes $y_1$, and the
terminal argument is removed against the vacuum. The integral is
exactly the stated vector formula. Pairing with $\Omega_1$ proves
\eqref{eq:rotation-ordered-density}.
A common real time translation puts any time-ordered supports in
the positive half-plane, without changing their differences or
their moment, proving the general formula.
For $N=1$, this formula and the vanishing mean field give $G_1=0$.

For the Gram identity, take the adjoint of the product defining
$\Psi(x;\lambda)$. It reverses the order, stars the arguments,
and multiplies by $s(\lambda)$. In the middle of the two products
the arguments $x_1^*$ and $z_1$ add to $x_1^*+z_1$; replace them
by equal halves. This gives
\eqref{eq:rotation-ordered-Gram}, including the stated vacuum cases.
The translation identity follows by shifting the first argument.

Conjugating \eqref{eq:rotation-ordered-density} replaces the plates
by $\bar\alpha_a\otimes\bar\lambda_a$. Testing in small ordered
product neighborhoods identifies
$\overline{G_N(\xi;\lambda)}=G_N(\xi;\bar\lambda)$ on real
$\xi\in\mathcal W^{N-1}$. Both sides, with $\xi$ conjugated on
the first, are holomorphic, and this domain is connected.
The identity theorem proves the claimed conjugation formula.
\end{proof}

\subsection{The two-point angular estimate}
\label{subsec:rotation-two-point-estimate}

The operator bound for $N$ point plates gives angular growth of
order $N$. At two points a positive spectral measure gives the
precise growth order two. Closedness will remove the unwanted
endpoint frequencies.

Fix a coordinate reconstruction $j$. For scalar
$u\in L^1\cap L^2(\R^3)$ supported in $x_j\ge0$, put
\[
 \widetilde u(E,p_\perp)=
       \int u(x)e^{-x_jE+ip_\perp\cdot x_\perp}\,\dd x.
\]
There are complex Radon measures $\varrho^{(j)}_{ab}$,
$a,b\in\{1,2,3\}$, with
\begin{equation}\label{eq:rotation-one-field-spectrum}
 \langle[\mathbb B_a(u)]_j,[\mathbb B_b(u')]_j\rangle
       =\int\overline{\widetilde u}\,\widetilde u'\,
                                    \dd\varrho^{(j)}_{ab}.
\end{equation}
They have a common positive Radon measure $\varrho_0$ such that
$\varrho^{(j)}_{ab}=r_{ab}\varrho_0$ and the matrix $r$ is
positive semidefinite and Hermitian almost everywhere.
In particular, for vector coefficients $c,c'$,
the measures $\varrho_{cc'}=\sum_{a,b}\bar c_ac'_b\varrho_{ab}$
satisfy
\begin{equation}\label{eq:rotation-spectral-CS}
 \int|G_1G_2|\,\dd|\varrho_{cc'}|
 \le\left(\int|G_1|^2\,\dd\varrho_{cc}\right)^{1/2}
       \left(\int|G_2|^2\,\dd\varrho_{c'c'}\right)^{1/2}.
\end{equation}
All these measures are carried by the joint spectral cone.

Here is the construction, including the fact that the measures
need only be locally finite. Convolution and
\exteq{6}{ErgodicComponents} give
\[
 [\mathbb B_a(u*\alpha)]_j
       =\widetilde u(H_j,\mathsf P)[\mathbb B_a(\alpha)]_j
\]
for $\alpha\in L^2$ supported in $x_j\ge0$.
This follows by applying the bounded field and quotient maps to
the $L^2$ Bochner integral
$u*\alpha=\int u(x)\tau_x\alpha\,\dd x$;
Fubini applies to its scalar spectral matrix elements since
$u\in L^1$.
Let $P_1(v)=(\pi(1+v^2))^{-1}$, whose Fourier transform with
either sign is $e^{-|p|}$, and choose
\[
 \alpha_0(x)=\mathbf1_{(0,1)}(x_j)\prod_{i\ne j}P_1(x_i),
 \qquad
 a_0(E,p_\perp)=\widetilde\alpha_0
       =\frac{1-e^{-E}}{E}e^{-|p_\perp|_1}>0.
\]
The quotient at $E=0$ is one. With
$\psi_a=[\mathbb B_a(\alpha_0)]_j$ and
$m_{ab}=\langle\psi_a,\mathsf E_j(\cdot)\psi_b\rangle$,
commutativity of convolution gives
\begin{equation}\label{eq:rotation-spectral-reference}
 a_0(H_j,\mathsf P)[\mathbb B_a(u)]_j
       =\widetilde u(H_j,\mathsf P)\psi_a.
\end{equation}
The finite matrix measure $m$ is positive. If $m_0=\sum_a m_{aa}$,
spectral Cauchy--Schwarz shows $m_{ab}\ll m_0$.
Its density matrix is positive Hermitian on one common full-measure
set: test its quadratic form on the countable rational complex
vectors and then use continuity. Define
$\varrho_{ab}=a_0^{-2}m_{ab}$ and $\varrho_0=a_0^{-2}m_0$.
The strictly positive continuous $a_0$ has a positive lower bound
on bounded sets, so these measures are Radon.
Pointwise matrix Cauchy--Schwarz followed by its integral version
gives \eqref{eq:rotation-spectral-CS}.

On $A_n=\{a_0>1/n\}$, which is bounded, one may divide
\eqref{eq:rotation-spectral-reference} by $a_0$ using bounded
spectral calculus. Thus
\[
 \big\langle\mathsf E_j(A_n)[\mathbb B_a(u)]_j,
             \mathsf E_j(A_n)[\mathbb B_b(u')]_j\big\rangle
       =\int_{A_n}\overline{\widetilde u}\,\widetilde u'
                                                  \,\dd\varrho_{ab}.
\]
As $A_n$ increase to the whole spectral space and
$\mathsf E_j(A_n)\to I$ strongly,
the diagonal identity follows by monotone convergence; its integral
is $\|[\mathbb B_a(u)]_j\|^2\le\|u\|_2^2$.
Equation \eqref{eq:rotation-spectral-CS} then makes the mixed
integral absolutely convergent and gives the unrestricted identity
by dominated convergence. Cone support follows from that of
$\mathsf E_j$.

Now use $j=1$, and take real profiles
$\chi,\chi'\in C_c^\infty(\R)$. Their transverse transforms in this
spectral convention are
$\widehat\chi_+(p_3)=\int\chi(x_3)e^{ip_3x_3}\,\dd x_3$.
For labels $\lambda=(\sigma,\chi)$ and
$\lambda'=(\sigma',\chi')$, set
\[
 \varrho_{\lambda\lambda'}=
 \sum_{a,b}\overline{e^\sigma_a}e^{\sigma'}_b
       \overline{\widehat\chi_+}\widehat{\chi'}_+\,\varrho^{(1)}_{ab}.
\]
In this formula $\widehat{\chi'}_+$ means the same positive-sign
transform of $\chi'$. These are Radon measures on $(E,p_2,p_3)$,
and $\varrho_{\lambda\lambda}\ge0$.
For $\alpha,\alpha'\in L^1\cap L^2(\R^2;\C)$ vanishing almost
everywhere on $y^0<0$, put
$\alpha^L(E,p_2)=\int\alpha(y)e^{-y^0E+iy^1p_2}\,\dd y$.
Equation \eqref{eq:rotation-one-field-spectrum} gives
\begin{equation}\label{eq:rotation-plate-spectrum}
 \langle[\mathbb B(\alpha\otimes\lambda)]_1,
          [\mathbb B(\alpha'\otimes\lambda')]_1\rangle
 =\int\overline{\alpha^L}(\alpha')^L\,
                                     \dd\varrho_{\lambda\lambda'}.
\end{equation}

The crucial estimate is
\begin{equation}\label{eq:rotation-plate-Laplace-bound}
 I_\lambda(t):=\int e^{-tE}\,\dd\varrho_{\lambda\lambda}
       \le\frac{36\|\chi\|_2^2}{\pi t^2}\quad(t>0).
\end{equation}
Indeed, for $c>0$ take
$\alpha_c(y)=c^{-2}\mathbf1_{(c,2c)}(y^0)P_1(y^1/c)$.
On $E\ge|p_2|$ its transform is bounded below by
$e^{-2cE-c|p_2|}\ge e^{-3cE}$.
Since $\|\alpha_c\|_2^2=(2\pi c^2)^{-1}$, the variance bound and
the quotient contraction in \eqref{eq:rotation-plate-spectrum} give
$I_\lambda(6c)\le\|\chi\|_2^2/(\pi c^2)$.
Choosing $c=t/6$ proves the estimate. Its mixed version is
\begin{equation}\label{eq:rotation-plate-mixed-Laplace}
 \int e^{-tE}\,\dd|\varrho_{\lambda\lambda'}|
       \le\sqrt{I_\lambda(t)I_{\lambda'}(t)}.
\end{equation}

Define, for $y^0>0$,
\[
 G^{[1]}_{\lambda\lambda'}(y)
       =s(\sigma)\int e^{-y^0E+iy^1p_2}\,
                                   \dd\varrho_{\lambda\lambda'}.
\]
It is smooth, since on $y^0\ge t>0$ every derivative of order $q$
is bounded in modulus by
$E^qe^{-tE}\le q!(2/t)^qe^{-tE/2}$ on the cone.
It computes the two-point function on time-ordered compact plane
supports. To see this, translate the supports to opposite sides
of $y^0=0$ and put $\alpha^\theta=\bar\alpha\circ r_1$.
Then
\[
 \E_\mu[\mathbb B(\alpha\otimes\lambda)
                 \mathbb B(\alpha'\otimes\lambda')]
 =s(\sigma)\int\overline{(\alpha^\theta)^L}(\alpha')^L\,
                                    \dd\varrho_{\lambda\lambda'}.
\]
Expanding the transforms gives the integral against
$G^{[1]}_{\lambda\lambda'}(y'-y)$. Fubini is justified by
\[
 \|\alpha\|_1\|\alpha'\|_1e^{-\delta E},
\]
where $\delta>0$ is the time gap, and
\eqref{eq:rotation-plate-mixed-Laplace}.
Thus this kernel agrees with $G_2$ on real positive-time differences.

For $r>0$ and $n_\zeta=(\cos\zeta,\sin\zeta)$, it has the
holomorphic continuation
\[
 h^{[1]}_{\lambda\lambda'}(\zeta)
   =s(\sigma)\int
       e^{-rE\cos\zeta+irp_2\sin\zeta}\,
                                \dd\varrho_{\lambda\lambda'}
       \quad(|\Re\zeta|<\pi/2).
\]
If $\zeta=t+i\beta$, the real part of the exponent is
$-r\cos t(E\cosh\beta+p_2\sinh\beta)$.
The cone bounds this above by
$-r\cos t\,e^{-|\beta|}E$.
On every compact substrip this gives an integrable exponential
majorant; dominated convergence and Morera's theorem give
holomorphy. Equations
\eqref{eq:rotation-plate-Laplace-bound}--\eqref{eq:rotation-plate-mixed-Laplace}
give
\begin{equation}\label{eq:rotation-two-point-strip}
 |h^{[1]}_{\lambda\lambda'}(t+i\beta)|
 \le\frac{36\|\chi\|_2\|\chi'\|_2}
          {\pi(r\cos t)^2}\,e^{2|\beta|}.
\end{equation}

\subsection{Closedness removes the unwanted two-point frequencies}
\label{subsec:rotation-two-point-frequencies}

Assume in addition that $\mu$ is closed and cubically invariant.
In particular, its two-point form is invariant under
$R_{\pi/2}$. We use only this quarter-turn consequence here.

For $\upsilon\in(\pi/2)\Z$ and $n_\upsilon\cdot y>0$, define
\[
 G^{(\upsilon)}_{\lambda\lambda'}(y)
   =e^{iw\upsilon}
       G^{[1]}_{\lambda\lambda'}(R_{-\upsilon}y),
       \qquad w=w(\sigma)+w(\sigma').
\]
By quarter-turn invariance and
\eqref{eq:rotation-plate-actions}, this is the two-point density
for supports ordered along $n_\upsilon$.
Two such densities agree on overlaps: test them on small product
neighborhoods ordered in both directions and use continuity.
The four half-planes cover $\R^2\setminus\{0\}$, so they define
one smooth kernel $G_{\lambda\lambda'}$ there.
It computes the two-point function for any two disjoint compact
plane supports. Indeed, partition them into sufficiently small
half-open squares. Each pair of squares is ordered along one of
the four coordinate directions, and there are only finitely many
pairs. Sum the ordered density identities.

For fixed $r>0$, the functions
\[
 h^{(\upsilon)}(\zeta)=e^{iw\upsilon}
       h^{[1]}_{\lambda\lambda'}(\zeta-\upsilon)
       \quad(|\Re\zeta-\upsilon|<\pi/2)
\]
agree on adjacent strips, by the identity theorem applied on
their real overlap. They give an entire function $h$ with
\begin{equation}\label{eq:rotation-two-point-entire-angle}
 h(\vartheta)=G_{\lambda\lambda'}(rn_\vartheta),\qquad
 h(\zeta+\pi/2)=i^w h(\zeta),\qquad
 |h(t+i\beta)|\le C_{\chi,\chi'}r^{-2}e^{2|\beta|},
\end{equation}
where $C_{\chi,\chi'}=72\|\chi\|_2\|\chi'\|_2/\pi$.
For the bound, choose a strip center within $\pi/4$ of $t$ and
apply \eqref{eq:rotation-two-point-strip}.

The entire function $e^{-iw\zeta}h(\zeta)$ has period $\pi/2$,
so it is $H(e^{4i\zeta})$ for a holomorphic $H$ on
$\C\setminus\{0\}$. The Laurent series of $H$ gives
frequencies $k=w+4m$, $m\in\Z$. Cauchy's formula on the circle
parametrized by $e^{4i(t+i\beta)}$, $0\le t\le\pi/2$,
bounds the coefficient $a^k(r)$ by
$C_{\chi,\chi'}r^{-2}e^{2|\beta|+k\beta}$.
Sending $\beta$ to the appropriate infinity kills $|k|>2$.
Consequently
\begin{equation}\label{eq:rotation-two-point-window}
 G_{\lambda\lambda'}(rn_\vartheta)
   =\sum_{\substack{|k|\le2\\k-w\in4\Z}}a^k(r)e^{ik\vartheta},
 \qquad |a^k(r)|\le C_{\chi,\chi'}r^{-2}.
\end{equation}
These coefficients are smooth in $r$ and real bilinear in the
profiles, by their Fourier integral formula.

Every pair of labels has only its expected frequency $k=w$,
except for the pairs $(+,+)$ and $(-,-)$, whose windows contain
both $2$ and $-2$. To eliminate the opposite frequency, use
closedness in its vector form $\mathbb B(\nabla u)=0$.
Take $\eta\in C_c^\infty(\R^2;\C)$ and
$\chi_0\in C_c^\infty(\R;\R)$. With $z=y^0+iy^1$ and
$\partial_z=(\partial_{y^0}-i\partial_{y^1})/2$,
\[
 \nabla(\eta(y)\chi_0(x_3))
   =(\partial_z\eta)\otimes(+,\chi_0)
       +(\partial_{\bar z}\eta)\otimes(-,\chi_0)
       +\eta\otimes(3,\chi'_0).
\]
Pair this gradient with a second plate whose plane support is
disjoint. Apply the density identities and integrate by parts.
Testing on small disjoint product neighborhoods gives
\begin{equation}\label{eq:rotation-two-point-Bianchi}
 \partial_zG_{(+,\chi_0)\lambda'}
       +\partial_{\bar z}G_{(-,\chi_0)\lambda'}
       +G_{(3,\chi'_0)\lambda'}=0
                       \quad\text{on }\R^2\setminus\{0\}.
\end{equation}
There is no diagonal contribution in this equation.
In polar coordinates,
\begin{align*}
 \partial_z(f(r)e^{ik\vartheta})
       &=\tfrac12e^{i(k-1)\vartheta}(f'(r)+kf(r)/r),\\
 \partial_{\bar z}(f(r)e^{ik\vartheta})
       &=\tfrac12e^{i(k+1)\vartheta}(f'(r)-kf(r)/r).
\end{align*}
For $\lambda'=(+,\chi')$, the frequency $-3$ in
\eqref{eq:rotation-two-point-Bianchi} occurs only in the derivative
of $a^{-2}_{(+,\chi_0)(+,\chi')}$.
It obeys $f'=2f/r$.
For $\lambda'=(-,\chi')$, the frequency $3$ likewise gives
$f'=2f/r$ for $a^2_{(-,\chi_0)(-,\chi')}$.
Each solution is $f(r)=f(1)r^2$. The bound
$|f(r)|\le C_{\chi_0,\chi'}r^{-2}$ for all $r>0$ forces
$f(1)=0$ by sending $r\to\infty$. Thus
\begin{equation}\label{eq:rotation-two-point-phase}
 G_{\lambda\lambda'}(R_ty)
       =e^{i(w(\sigma)+w(\sigma'))t}G_{\lambda\lambda'}(y).
\end{equation}
The plate action cancels this phase in the two-point density
formula, proving rotational covariance for disjoint compact
plane supports.

\subsection{Extending identities across coincident plane supports}
\label{subsec:rotation-contact}

The preceding density argument does not yet address coincident
plane supports. We use the following $L^2$ argument,
which will also be applied at every higher order.

\begin{proposition}\label{prop:rotation-no-contact}
Let $U\subset\R^2$ be open, and write $L^2_U$ for the vector
fields supported in $U\times\R$. Let $D:(L^2_U)^N\to\C$ be a
bounded multilinear form. Suppose it vanishes when the entries
are supported over pairwise disjoint compact subsets of $U$,
and whenever one entry is $\nabla u$ with
$u\in C_c^\infty(U\times\R)$. Then $D=0$.
\end{proposition}

\begin{proof}
We argue simultaneously for all open $U$ by induction on $N$.
At $N=1$, fields with compact plane support are dense, so the
first assumption suffices.

At $N=2$, represent the bilinear form as
$D(f_1,f_2)=\int f_1(x)\cdot(Tf_2)(x)\,\dd x$ for a bounded
operator $T$ on $L^2_U$. The pairing here is bilinear.
For disjoint compact $K,K'\subset U$ the first assumption gives
$\mathbf1_KT\mathbf1_{K'}=0$.
Inner regularity and strong convergence of the corresponding
multiplication projections extend this to disjoint Borel sets.
Using $A$ and $U\setminus A$ on both sides gives
$T\mathbf1_A=\mathbf1_AT$ for every Borel $A$.
Uniform approximation by simple functions then gives $Th=hT$
for every bounded plane multiplier $h$. Hence
\[
 D(h\nabla u,f_2)=D(\nabla u,hf_2)=0.
\]

Fix a compact $A\subset U$, choose $u_0\in C_c^\infty(U)$ equal
to one near $A$, and let $\chi_0\in C_c^\infty(\R)$.
On $A\times\R$, the gradients of
$u_0\chi_0$, $u_0y^0\chi_0$, and $u_0y^1\chi_0$ are respectively
\[
 \chi_0'e_3,\qquad
 \chi_0e_1+y^0\chi_0'e_3,\qquad
 \chi_0e_2+y^1\chi_0'e_3.
\]
The preceding identity with multiplier $h\mathbf1_A$ first
annihilates $h\mathbf1_A\chi'_0e_3$.
Applying it with $hy^0\mathbf1_A$ and $hy^1\mathbf1_A$ then
annihilates $h\mathbf1_A\chi_0e_1$ and
$h\mathbf1_A\chi_0e_2$. These multipliers are bounded since $A$
is compact. Such plane-transverse products span a dense subspace.
For the transverse derivative family, its orthogonal complement
in $L^2(\R)$ has distributional derivative zero, hence is a
constant and therefore zero. Thus these three vector families
are total in $L^2_U$, and $D=0$.

Suppose $N\ge3$ and the assertion is proved at every smaller order.
Two consequences of the induction hypothesis separate unequal
plane cells. First, if one entry is supported over a compact
$K\subset U$ and the others over $U\setminus K$, fix that entry.
The form on the remaining $N-1$ slots satisfies both assumptions
on the open set $U\setminus K$ and is zero.
Second, if the slots divide into nonempty groups $I,J$, with
the $J$ entries supported over a compact $K_J$ and the $I$
entries over its complement, fix the $J$ entries.
The form on the $I$ slots inherits the gradient assumption.
When those slots have pairwise disjoint compact supports, the
first consequence, applied to any one of them, gives zero.
The induction hypothesis on $U\setminus K_J$ therefore makes
this form zero as well.

By density put all entries over one compact $K_0\subset U$.
Partition the plane into half-open squares $Q$ of side $d$ and
expand $D$ by multilinearity. A tuple of squares which are not
all equal has zero contribution. To justify compact separation,
replace each square by its concentric closed square of side
$d-2\epsilon$. Separate the slots whose square equals the first
from the remaining slots, use the second consequence with
$K_J=K_0$ intersected with the union of their closed squares,
and send $\epsilon\downarrow0$ in $L^2$.
Only the equal-square terms remain, and Cauchy--Schwarz gives
\[
 |D(f_1,\ldots,f_N)|
 \le\|D\|\max_Q\|\mathbf1_Qf_3\|_2
       \|f_1\|_2\|f_2\|_2\prod_{i=4}^N\|f_i\|_2.
\]
Indeed, sum the first two square norms by Cauchy--Schwarz and
bound the others by their whole norms. The maximum tends to
zero as $d\downarrow0$: its square is the integral over an
area-$d^2$ cell of the fixed integrable function
$y\mapsto\int|f_3(y,x_3)|^2\,\dd x_3$.
Absolute continuity of the integral proves the assertion.
\end{proof}

Apply this proposition with $N=2$ to
\[
 D(b,b')=\E_\mu[\mathbb B(R_tb)\mathbb B(R_tb')]
              -\E_\mu[\mathbb B(b)\mathbb B(b')],
\]
where $R_tb$ means $R_t\cdot b$. Gaussian domination makes it a
bounded bilinear form. Finite sums of plates with smooth real
transverse profiles are dense in the fields over any fixed
compact plane set: simple plane-transverse tensors are dense
in $L^2$, and $e^+,e^-,e^3$ span $\C^3$.
Thus \eqref{eq:rotation-two-point-phase} supplies the disjoint-support
assumption for general fields. Closedness supplies the gradient
assumption, since rotations carry gradients to gradients.
Consequently the two-point function is invariant under every
$R_t$, on all $L^2$ tests.

\subsection{The angular expansion at every order}
\label{subsec:rotation-angular-expansion}

Fix $N\ge2$ and a tuple of labels $\lambda$, with total weight
$w=w(\lambda)$. Put $\mathcal W_\upsilon=R_\upsilon\mathcal W$
for $\upsilon\in(\pi/2)\Z$. For a permutation $\mathfrak s$,
let
\[
 \mathcal Z(\mathfrak s,\upsilon)=
 \{Z=(z_1,\ldots,z_N):
       z_{\mathfrak s(a+1)}-z_{\mathfrak s(a)}
                    \in\mathcal W_\upsilon,\ 1\le a<N\}.
\]
On this open convex set define
\[
 G_\upsilon^{\mathfrak s}(Z;\lambda)=e^{iw\upsilon}
 G_N\left(
       (R_{-\upsilon}(z_{\mathfrak s(a+1)}-z_{\mathfrak s(a)}))_{a=1}^{N-1};
       \mathfrak s\lambda\right).
\]
Here $(\mathfrak s\lambda)_a=\lambda_{\mathfrak s(a)}$.
Moment symmetry, quarter-turn invariance, and
\eqref{eq:rotation-ordered-density} identify this as the moment
density on compact product supports in its real part.
Two such holomorphic functions agree on their intersection.
Indeed, the intersection is convex and is preserved by taking
real parts, so a nonempty intersection contains an open set of
real configurations. Testing both densities there identifies
them by continuity, and the holomorphic identity theorem
identifies them on the full intersection.
They therefore define one holomorphic function
$\mathcal G_N(Z;\lambda)$ on
$\mathcal Z_N=\bigcup_{\mathfrak s,\upsilon}\mathcal Z(\mathfrak s,\upsilon)$.

The real set $\mathcal O_N=\mathcal Z_N\cap(\R^2)^N$ consists
of configurations whose first coordinates are all distinct
or whose second coordinates are all distinct.
Its complement is the finite union of the codimension-two
subspaces
$\{y_a^0=y_b^0,\ y_k^1=y_l^1\}$.
Thus $\mathcal O_N$ is open, dense, and connected.
The functions $\mathcal G_N$ are symmetric under simultaneous
permutation of points and labels, invariant under common
translations, and satisfy
\[
 \mathcal G_N(R_{\pi/2}Z;\lambda)=i^w\mathcal G_N(Z;\lambda),
 \qquad
 \overline{\mathcal G_N(Z;\lambda)}
       =\mathcal G_N(\bar Z;\bar\lambda).
\]
These identities follow from the defining densities and
the corresponding identities for $G_N$.

For $\beta\in\R$ write $\Lambda_\beta=R_{i\beta}$.
In the coordinates $\xi^\pm=\xi^0\pm i\xi^1$,
$\mathcal W=\{\Re\xi^+>0,\Re\xi^->0\}$ and
$\Lambda_\beta\xi^\pm=e^{\mp\beta}\xi^\pm$.
Thus every tube $\mathcal Z(\mathfrak s,\upsilon)$ is preserved
by $\Lambda_\beta$. Moreover
$d_{\mathcal W}(\Lambda_\beta\xi)\ge
e^{-|\beta|}d_{\mathcal W}(\xi)$ on the closed tube.

\begin{proposition}\label{prop:rotation-finite-angular-expansion}
There are holomorphic coefficients $c_m(Z;\lambda)$ on
$\mathcal Z_N$, vanishing unless $|w+4m|\le N$, such that
\begin{equation}\label{eq:rotation-finite-angular-expansion}
 \mathcal G_N(\Lambda_\beta Z;\lambda)
       =\sum_m c_m(Z;\lambda)e^{-(w+4m)\beta}
                         \quad(\beta\in\R).
\end{equation}
They are unique, symmetric under simultaneous permutation,
and invariant under common translations. They satisfy
\begin{align}
 c_m(\Lambda_\beta Z;\lambda)
       &=e^{-(w+4m)\beta}c_m(Z;\lambda),
                           \label{eq:rotation-angular-coefficient-identities}\\
 c_m(R_{\pi/2}Z;\lambda)&=i^wc_m(Z;\lambda),\qquad
 \sum_m c_m(Z;\lambda)=\mathcal G_N(Z;\lambda),\notag\\
 \overline{c_m(X;\lambda)}&=c_{-m}(X;\bar\lambda)
                           \quad(X\in\mathcal O_N).\notag
\end{align}
\end{proposition}

\begin{proof}
Fix $0<\vartheta_0<\pi/4$ and first consider real configurations
whose consecutive differences have positive lengths and angles
$|\vartheta_a|<\vartheta_0$ from $n_0$.
For $\zeta=t+i\beta$,
\[
 \Re(R_{\zeta-\upsilon}\xi_a)^\pm
       =|\xi_a|e^{\mp\beta}\cos(\vartheta_a+t-\upsilon).
\]
The strips $|t-\upsilon|<\pi/2-\vartheta_0$ cover the complex
$\zeta$-plane. Thus $h_X(\zeta)=\mathcal G_N(R_\zeta X;\lambda)$
is entire. Choosing $\upsilon$ within $\pi/4$ of $t$ bounds every
gap distance below by
$r_Xe^{-|\beta|}\cos(\vartheta_0+\pi/4)$, where
$r_X=\min_a|\xi_a|$.
Equation \eqref{eq:rotation-ordered-kernel-bound} gives
$|h_X(t+i\beta)|\le C_{X,\lambda}e^{N|\beta|}$.
Its quarter-period is $h_X(\zeta+\pi/2)=i^wh_X(\zeta)$.
The Laurent coefficient argument used at two points, with
growth exponent $N$, gives the finite frequencies
$k=w+4m$, $|k|\le N$.

To extend this statement to a whole tube, let
$\mathcal M=\{m:|w+4m|\le N\}$ and
$P(u)=\prod_{m\in\mathcal M}(u+w+4m)$.
The functions
$Z\mapsto P(\partial_\beta)G_\upsilon^{\mathfrak s}
(\Lambda_\beta Z;\lambda)$ are holomorphic for each real $\beta$.
They vanish on the nonempty open real subset obtained by rotating
and permuting the configurations just used.
The identity theorem gives zero on the entire convex tube.
The resulting constant-coefficient ordinary differential equation
has distinct characteristic roots $-(w+4m)$.
Its coefficients in the exponential expansion are fixed linear
combinations of its derivatives at $\beta=0$, and hence are
holomorphic in $Z$.
On overlapping tubes the exponentials are linearly independent,
so the coefficients agree and define $c_m$ on $\mathcal Z_N$.
This also proves uniqueness. Comparing coefficients after a
permutation, common translation, or addition of the parameter
$\beta$ gives the stated identities. For the conjugation identity
use $\overline{\Lambda_\beta X}=\Lambda_{-\beta}X$ and
$w(\bar\lambda)=-w$.
\end{proof}

The coefficients have exact real rotation phases whenever both
configurations are in $\mathcal O_N$:
\begin{equation}\label{eq:rotation-angular-coefficient-phase}
 c_m(R_tX;\lambda)=e^{i(w+4m)t}c_m(X;\lambda).
\end{equation}
Locally in one tube, the allowed angles form an interval:
each ordered difference has a positive projection precisely
when $\cos(\vartheta_a+t)>0$.
The identity theorem in $t+i\beta$ extends
\eqref{eq:rotation-finite-angular-expansion} from the imaginary
axis to that interval. Comparing the exponential coefficients
gives \eqref{eq:rotation-angular-coefficient-phase} locally.

To pass through changes of ordering, first take distinct points
with no two distinct chords perpendicular. Such configurations
are a dense open set, the complement of finitely many nonzero
polynomial zero sets. Every rotation of one of them belongs to
$\mathcal O_N$: simultaneous coincidences of first and second
coordinates would give perpendicular chords, or identical points.
The local phase identity therefore continues along its whole
rotation orbit. Approximation proves it whenever $X,R_tX$
are in $\mathcal O_N$.

For an arbitrary configuration $X$ of distinct plane points,
choose $t$ so that all first coordinates of $R_tX$ are distinct.
Only finitely many angles modulo $\pi$ are excluded.
Define
\[
 \widetilde c_m(X;\lambda)
       =e^{-i(w+4m)t}c_m(R_tX;\lambda).
\]
The phase identity makes this independent of the choice.
Locally one may keep $t$ fixed, so it is real analytic on the
connected space $(\R^2)^N_{\ne}$ of distinct configurations.
It inherits permutation symmetry, translation invariance, and
\[
 \widetilde c_m(R_tX;\lambda)
       =e^{i(w+4m)t}\widetilde c_m(X;\lambda).
\]
Put $\widetilde G_N=\sum_m\widetilde c_m$.
It is real analytic and agrees with $\mathcal G_N$ on $\mathcal O_N$.

Define $\mathfrak a(0)=\mathfrak a(1)=0$, and, for $N\ge2$,
let $\mathfrak a(N)$ be the largest $|m|$ for which $c_m$ is
not identically zero for some tuple of $N$ labels; set it to zero
if all coefficients vanish. Since $|w|\le N$ and $|w+4m|\le N$,
\begin{equation}\label{eq:rotation-angular-linear-bound}
 \mathfrak a(N)\le N/2.
\end{equation}
A nonzero coefficient can be detected on a strictly time-ordered
real configuration after a quarter turn and permutation:
the identity theorem on each tube and the coefficient identities
above give this reduction. Conjugation then shows that if
$\mathfrak a(N)>0$, both indices
$\mathfrak a(N)$ and $-\mathfrak a(N)$ occur.

\subsection{Reflection bounds for the angular coefficients}
\label{subsec:rotation-angular-reflections}

We next prove the convexity estimate
\begin{equation}\label{eq:rotation-angular-convexity}
 2\mathfrak a(a+b)\le\mathfrak a(2a)+\mathfrak a(2b)
                          \quad(a,b\ge0).
\end{equation}
Let $N=a+b$ and $M=\mathfrak a(N)>0$.
Choose a strictly time-ordered $X=(y_1,\ldots,y_N)$ and labels
with $c_{-M}(X;\lambda)\ne0$. Translate its time cut so that the
first $a$ points have negative times and the remaining $b$
have positive times. Reflect the first group and reverse its
order to obtain $x$, with reversed labels $\lambda^x$;
call the second group $z$, with labels $\lambda^z$.
Their weights satisfy $w=w_x+w_z$.
The Gram identity gives
\[
 \langle\Psi(\Lambda_\beta x;\lambda^x),
             \Psi(\Lambda_\beta z;\lambda^z)\rangle
       =s(\lambda^x)\mathcal G_N(\Lambda_\beta X;\lambda).
\]
An empty group gives the vacuum.
The squared norm of the first vector is, up to the sign
$s(\lambda^x)$, the kernel of its reflected double:
the configuration $(r_1x_a,\ldots,r_1x_1,x_1,\ldots,x_a)$
with the same labels in reverse then forward order.
Its weight is $2w_x$. The other squared norm has weight $2w_z$.
These are direct instances of
\eqref{eq:rotation-ordered-Gram}; the reflection-conjugation
keeps the labels.

Apply Cauchy--Schwarz to this matrix element and expand the
three kernels by \eqref{eq:rotation-finite-angular-expansion}.
Multiplication by $e^{2w\beta}$ cancels their weights and gives
\[
 \left|\sum_m c_m(X;\lambda)e^{-4m\beta}\right|^2
 \le
 \left|\sum_{|m|\le\mathfrak a(2a)}d_me^{-4m\beta}\right|
 \left|\sum_{|m|\le\mathfrak a(2b)}d'_me^{-4m\beta}\right|.
\]
The coefficients $d_m,d'_m$ belong to the two specified doubles.
A vacuum factor is one. As $\beta\to+\infty$, the left side
has a nonzero leading term of rate $8M$, and the right side is
bounded by a constant times
$e^{4(\mathfrak a(2a)+\mathfrak a(2b))\beta}$.
This proves \eqref{eq:rotation-angular-convexity}.
The cases $M=0$ or $N\le1$ already satisfy it.

There is also a positive kernel for the highest coefficient.
For $a\ge1$, set $M=\mathfrak a(2a)$. On labelled configurations
of $a$ distinct points in $y^0>0$, define
\begin{equation}\label{eq:rotation-top-coefficient-kernel}
 K_a((x;\lambda),(z;\lambda'))=
       s(\lambda)\widetilde c_{-M}^{(2a)}
                        (r_1x,z;\lambda,\lambda').
\end{equation}
The superscript indicates the number of points.
This kernel is Hermitian and positive semidefinite.
First suppose the times in each configuration are distinct and
order them increasingly. By the Gram formula,
\[
 \langle\Psi(\Lambda_\beta x;\lambda),
                 \Psi(\Lambda_\beta z;\lambda')\rangle
 =s(\lambda)\sum_{|m|\le M}
       \widetilde c_m^{(2a)}(r_1x,z;\lambda,\lambda')
       e^{-(w(\lambda)+w(\lambda')+4m)\beta}.
\]
Reversing the reflected points is harmless by labelled symmetry.
Multiply this Gram kernel by the positive factors
$e^{w(\lambda)\beta}e^{w(\lambda')\beta}e^{-4M\beta}$.
Its pointwise limit as $\beta\to+\infty$ is $K_a$.
Positive semidefiniteness passes to this limit on every finite
matrix. Approximating configurations with tied times by ones
with distinct times proves the same assertion for all the
stated configurations.

The exact phase of $\widetilde c_{-M}$ transports this consequence
to any line $L$ in the plane. If $x,z$ lie on one open side of
$L$ and $r_L$ is its reflection, translation to a point of $L$
and rotation of its normal to $n_0$ express
$\widetilde c_{-M}(r_Lx,z)$ as $K_a$ times a unit-modulus factor.
Taking absolute values of the positive-kernel
Cauchy--Schwarz inequality gives
\begin{multline}\label{eq:rotation-top-coefficient-CS}
 |\widetilde c_{-M}^{(2a)}(r_Lx,z;\lambda,\lambda')|^2\\
 \le
 |\widetilde c_{-M}^{(2a)}(r_Lx,x;\lambda,\lambda)|
 |\widetilde c_{-M}^{(2a)}(r_Lz,z;\lambda',\lambda')|.
\end{multline}
In particular a zero reflected diagonal forces every such
off-diagonal coefficient to vanish.

\subsection{Regular polygons and the reduction to two points}
\label{subsec:rotation-polygons}

We prove the arithmetic restriction
\begin{equation}\label{eq:rotation-polygon-divisibility}
 a\mid 2\mathfrak a(2a)\quad(a\ge1).
\end{equation}
Fix $a$ and put $M=\mathfrak a(2a)$,
$c_*=\widetilde c_{-M}^{(2a)}$.
The coefficients depend only on the labelled set of points.
If a rotation through $t$ permutes such a set without changing
its labels, its exact phase gives
$(1-e^{i(w-4M)t})c_*=0$.
For a regular $2a$-gon with the same label $(\sigma,\chi)$ at
each vertex, use $t=\pi/a$ and $w=2aw(\sigma)$.
The label weight cancels, leaving $e^{-4\pi iM/a}$.
Thus $c_*$ vanishes on this polygon whenever $a$ fails to divide $2M$.

That selection rule concerns a special configuration. The
positive kernel propagates its zero to arbitrary configurations
by the following halving argument. For a labelled set $S$ of
at most $2a$ distinct points, let $\mathcal N(S)$ mean that
$c_*$ vanishes on every $2a$-point labelled configuration
containing $S$.
Let $S'$ have $j'\le a$ points on one open side $L_+$ of a line
$L$, and suppose $S\subset r_LS'\cup S'$.
Then
\begin{equation}\label{eq:rotation-halving}
 \mathcal N(S)\ \Longrightarrow\
       \mathcal N(S')\ \text{and}\ \mathcal N(r_LS').
\end{equation}
To see this, add $a-j'$ arbitrary points $R$ in $L_+$ to $S'$,
forming $x=S'\cup R$. The reflected double $r_Lx\cup x$
contains $S$, so its diagonal coefficient is zero.
Equation \eqref{eq:rotation-top-coefficient-CS} implies
$c_*(r_Lx\cup z)=0$ for every $a$-point labelled set $z$ in $L_+$.
Thus, with $r_LS'$ fixed, the remaining $a-j'$ points on
$L_-$ and $a$ points on $L_+$ may vary in an open set.

For each choice of their labels, the space of their distinct
positions, avoiding the fixed points, is the complement of
finitely many codimension-two affine subspaces and is connected.
The coefficient is real analytic there. Its zero on the
nonempty open set just obtained extends to all positions.
Every $2a$-point configuration containing $r_LS'$ is obtained
this way. This proves $\mathcal N(r_LS')$.
Interchange the sides of $L$ and repeat the argument to obtain
$\mathcal N(S')$.

Suppose now $a\nmid2M$. Then $M>0$ and $a\ge3$.
Fix an arbitrary label and any regular $2a$-gon with vertices
$P_k$, $k\in\Z/(2a)\Z$, in cyclic order.
Let $\mathcal N_j$ assert $\mathcal N(S)$ for every arc of $j$
consecutive vertices, with this label, in every such polygon.
Selection gives $\mathcal N_{2a}$.
For $j\ge2$, put $j'=\lceil j/2\rceil\le a$ and take an arc
$S'=\{P_{k_0+1},\ldots,P_{k_0+j'}\}$.
Reflect in the perpendicular bisector of
$[P_{k_0},P_{k_0+1}]$.
The $a$ vertices $P_{k_0+1},\ldots,P_{k_0+a}$ lie on one
open side: their angles from the edge-midpoint radius are
$(2k-1)\pi/(2a)\in(0,\pi)$.
Their images satisfy $r_LP_{k_0+k}=P_{k_0+1-k}$.
Consequently $r_LS'\cup S'$ is an arc of $2j'$ consecutive
vertices and contains a $j$-vertex arc.
Equation \eqref{eq:rotation-halving} gives
$\mathcal N_j\Rightarrow\mathcal N_{j'}$.

Iteration reaches $j=1$. Any plane point is a vertex of some
regular $2a$-gon, so the coefficient vanishes whenever a
configuration has even one point with the fixed label.
Since the label was arbitrary, $c_*$ vanishes identically for
all label tuples. This contradicts the occurrence of the
index $-M$, proving \eqref{eq:rotation-polygon-divisibility}.

Now set $\nu(a)=\mathfrak a(2a)$. It is a nonnegative integer
sequence, with $\nu(0)=0$, $\nu(a)\le a$, and
$2\nu(a)\le\nu(a-1)+\nu(a+1)$ by
\eqref{eq:rotation-angular-convexity}.
Its increments $d_a=\nu(a)-\nu(a-1)$ are nondecreasing integers.
The linear bound keeps them bounded above, so they are eventually
constant, say $d$. Thus $\nu(a)=da-c$ for large $a$, where
$c=\sum_k(d-d_k)\ge0$ is an integer.
Divisibility \eqref{eq:rotation-polygon-divisibility} makes every
large $a$ divide $2c$, forcing $c=0$.
All increments equal $d=\nu(1)$, and
\[
 \mathfrak a(2a)=a\mathfrak a(2),\qquad
 2\mathfrak a(N)\le N\mathfrak a(2).
\]
For the second inequality split $N=\lfloor N/2\rfloor+
\lceil N/2\rceil$ in \eqref{eq:rotation-angular-convexity}.

The two-point result gives $\mathfrak a(2)=0$.
For smooth transverse profiles this follows from
\eqref{eq:rotation-two-point-phase} and the ordered density formula:
near any strictly ordered two-point configuration, only the
frequency $w$ occurs. Linear independence of the finite angular
exponentials makes $c_m=0$ there for $m\ne0$, and the identity
theorem on each tube gives zero throughout.
For arbitrary $L^2$ profiles, the all-$L^2$ two-point covariance
proved above gives the same local density identity.
Hence
\begin{equation}\label{eq:rotation-no-extra-angular-coefficients}
 \mathfrak a(N)=0\quad\text{for every }N,\qquad
 \widetilde G_N(R_tX;\lambda)=e^{iw(\lambda)t}
                                  \widetilde G_N(X;\lambda).
\end{equation}

\subsection{Moment densities on all distinct configurations}
\label{subsec:rotation-all-distinct-density}

The analytic continuation $\widetilde G_N$ must still be
identified as the moment density at distinct configurations
outside $\mathcal O_N$. This is a separate step; continuation
of a density does not by itself exclude an additional singular
part on that smaller set.

We first record the analytic displacement used for this step.
If polynomial field observables $\Phi'$ and $\Phi''$ have test
supports in $x_1<c-\delta$ and $x_1>c+\delta$, respectively,
then
\[
 s\longmapsto\E_\mu[\Phi'U(se_2)\Phi'']
\]
extends holomorphically to $|\Im s|<2\delta$.
Translate $c$ to zero and put
$\psi'=[U_{-\delta}\Theta_1\Phi']_1$,
$\psi''=[U_{-\delta}\Phi'']_1$.
The quotient and semigroup identities express the real
correlation as
$\langle\psi',T((2\delta,s))\psi''\rangle$.
This is holomorphic on the stated strip by the cone.

For compact, pairwise disjoint plane supports we claim
\begin{equation}\label{eq:rotation-distinct-density}
 \E_\mu\prod_a\mathbb B(\alpha_a\otimes\lambda_a)
 =\int\prod_a\alpha_a(y_a)\widetilde G_N(Y;\lambda)\,\dd Y.
\end{equation}
The integral is absolutely convergent: the product support
is a compact subset of $(\R^2)^N_{\ne}$, the kernel is bounded
there, and each $\alpha_a\in L^2$ is integrable on its compact
support.
Call the difference of the two sides $D(\alpha_1,\ldots,\alpha_N)$
on these tuples.

We show local vanishing near each distinct configuration $X$
by induction on the number of equal-second-coordinate pairs.
When there are none, $X\in\mathcal O_N$ and the ordered tube
density gives the identity in small product discs.
Otherwise take a pair with equal second coordinates and
different times. Choose a time cut between them, different
from all point times, and split the points into groups $I,J$
below and above the cut.
Move the $J$ group by $se_2$. The points remain distinct for
every real $s$, since the groups are separated in time.
A sufficiently small positive $s_0$ removes every
second-coordinate coincidence across the cut and creates none:
take it smaller than every nonzero absolute difference of the
initial second coordinates. If all these coordinates agree, any
sufficiently small positive $s_0$ works.
There are fewer coincident pairs at $X(s_0)$.

The induction hypothesis gives a radius $2r$ near $X(s_0)$
on which $D$ vanishes. Reduce $r$ so that the original discs
of radius $r$ are disjoint and preserve the time gap.
Fix tests in these original discs, and translate those in
group $J$. The resulting difference $F(s)$ is real analytic
on $\R$.
For its moment term this follows from the separated-cluster
strip argument above. For its density term, change variables
to integrate
$\widetilde G_N(Y+sv_J;\lambda)$ against the fixed product
of tests, where $(v_J)_a=e_2$ for $a\in J$ and zero otherwise.
Over any bounded real $s$-interval these configurations form
a compact set of distinct points. Finite local holomorphic
extensions of the real analytic kernel give a common radius
$\epsilon>0$ and a bound $C$ there. Cauchy's estimates bound
the integrated derivative of order $q$ by
$Cq!\epsilon^{-q}\prod_a\|\alpha_a\|_1$.
Taylor's formula therefore gives real analyticity of this
integral.

For $|s-s_0|<r$ the translated tests remain in the discs of
radius $2r$ at $X(s_0)$, so $F(s)=0$.
The real analytic identity theorem gives $F(0)=0$.
This proves local vanishing near $X$ and completes the induction.
For general disjoint compact supports, finitely many of these
product-disc neighborhoods cover their product.
Partition each support into sufficiently small half-open
squares. Every product of nonzero pieces lies in one of the
larger vanishing neighborhoods: choose a point of the
product support in it and make each square diameter smaller
than the available half-radius.
Multilinearity and finite summation prove
\eqref{eq:rotation-distinct-density}.

\subsection{Euclidean covariance and all reflection planes}
\label{subsec:rotation-conclusion}

\begin{theorem}\label{thm:rotation-covariance}
Let $\mu$ be a Gaussian-dominated law of closed tempered two-forms.
Suppose it is invariant under translations and the cubic group,
and reflection positive at the nine planes
$x_j=0$ and $x_j\pm x_i=0$, $i<j$.
Then it is invariant under all Euclidean motions and reflection
positive at every plane, on the full $L^2$ space of either open side.
\end{theorem}

\begin{proof}
Proposition~\extref{3.1}{ErgodicComponents} supplies the sharp coordinate
cones. For the plane $(1,2)$, fix $t$ and define the bounded
$N$-linear form
\[
 D(b_1,\ldots,b_N)=
 \E_\mu\prod_a\mathbb B(R_t\cdot b_a)
       -\E_\mu\prod_a\mathbb B(b_a).
\]
Its boundedness follows from Proposition~\extref{2.1}{ContinuumLimits}.
For plates over pairwise disjoint compact plane supports,
\eqref{eq:rotation-distinct-density} and
\eqref{eq:rotation-no-extra-angular-coefficients} prove $D=0$:
the factor $e^{-iw(\lambda)t}$ from the plate action cancels
the factor $e^{iw(\lambda)t}$ from the kernel.
Density of finite sums of plates gives the same assertion for
arbitrary $L^2$ fields over those compact sets.
The form also vanishes on gradients, by closedness and the
fact that rotations carry gradients to gradients.
\Cref{prop:rotation-no-contact} therefore gives $D=0$
on all $L^2$ entries, for every $N$.

The entire generating function of each finite real Schwartz
test family is determined by these moments. Thus its
characteristic function is unchanged by $R_t$, and Fourier
uniqueness on the evaluation sigma algebra gives
$(R_t)_*\mu=\mu$.
Permuting coordinates applies the same argument in the other
coordinate planes: the nine mirrors give their cones, and
cubic invariance gives their quarter turns.
These plane rotations generate $SO(3)$.
A coordinate reflection already preserves $\mu$, so the full
$O(3)$ acts invariantly. Translation invariance gives the
Euclidean group.

Let $g$ carry $x_1=0$ to any chosen plane and its positive side
to the chosen open side. The composition operator of $g$ is
unitary on $L^2(\mu)$, maps these side spaces onto each other,
and conjugates the axis reflection to the chosen reflection.
Transporting its nonnegative reflection form proves the final
assertion.
\end{proof}

\Cref{thm:rotation-covariance} applies separately to the
component laws of Sections~\extref{5--6}{ErgodicComponents}.
They retain Gaussian domination, closedness, the cubic group,
and the nine mirror forms. Together with their all-direction
clustering, it supplies the hypotheses for the tensor OS
verification.



\printbibliography
\end{document}
